\documentclass[11pt,letterpaper]{article}
\usepackage[T1]{fontenc}
\usepackage{lmodern,microtype}
\usepackage[margin=1in]{geometry}
\usepackage{amsmath,amssymb,amsthm,mathtools,booktabs,array,longtable}
\usepackage{xcolor,enumitem,fancyhdr}
\usepackage[colorlinks=true,linkcolor=blue!45!black,citecolor=blue!45!black,urlcolor=blue!45!black,breaklinks=true]{hyperref}
\pdfmapfile{+lm.map}\pdfmapfile{+cm.map}\pdfmapfile{+symbols.map}
\pdfinfoomitdate=1\pdfsuppressptexinfo=15
\pdftrailerid{<5371756172655461626C65733230323600><5371756172655461626C65733230323600>}
\hypersetup{pdftitle={All Fixed Order Asymptotics for Square Contingency Tables},pdfauthor={Report 236},pdfsubject={OEIS A110058 geometric saddle expansion and controlled inverse}}
\pagestyle{fancy}\fancyhf{}\fancyhead[L]{\small Square contingency tables}\fancyhead[R]{\small Report 236}\fancyfoot[C]{\thepage}
\renewcommand{\headrulewidth}{0pt}\setlength{\headheight}{14pt}
\setlength{\parskip}{4pt}\setlength{\emergencystretch}{2em}
\setlist{itemsep=2pt,topsep=4pt}\numberwithin{equation}{section}
\newtheorem{theorem}{Theorem}[section]\newtheorem{lemma}[theorem]{Lemma}\newtheorem{proposition}[theorem]{Proposition}\newtheorem{corollary}[theorem]{Corollary}
\theoremstyle{definition}\newtheorem{definition}[theorem]{Definition}
\theoremstyle{remark}\newtheorem{remark}[theorem]{Remark}
% Added by the ProveIt write phase (batch 108, 6 October 2026): dated notes
% marked [write] and a macro for file, label and declaration names.
\newcommand{\file}[1]{\nolinkurl{#1}}
\expandafter\def\expandafter\UrlBreaks\expandafter{\UrlBreaks\do\-\do\_\do\.}
\newenvironment{writenote}[1][]{\par\smallskip\begingroup\small
 \noindent\textbf{[write]}\if\relax\detokenize{#1}\relax\else~\textbf{#1.}\fi\ }%
 {\par\endgroup\smallskip}
\newcommand{\E}{\mathbb E}\newcommand{\PP}{\mathbb P}\newcommand{\R}{\mathbb R}\newcommand{\N}{\mathbb N}\newcommand{\ii}{\mathrm i}\newcommand{\1}{\mathbf 1}\newcommand{\Cov}{\operatorname{Cov}}\newcommand{\Var}{\operatorname{Var}}\newcommand{\eps}{\varepsilon}\newcommand{\cost}{\operatorname{cost}}\newcommand{\Part}{\mathcal P}\newcommand{\dd}{\mathrm d}\newcommand{\ceil}[1]{\left\lceil #1\right\rceil}
\begin{document}
\begin{center}
{\large Report 236}\par\vspace{8pt}
{\LARGE\bfseries All Fixed Order Asymptotics\\[5pt] for Square Contingency Tables}\par\vspace{9pt}
{\large The density one model and its controlled inverse}\par\vspace{7pt}
5 October 2026
\end{center}
\begin{abstract}
Let $a_n$ count labelled $n\times n$ nonnegative integer matrices with every row and column sum equal to $n$, the sequence OEIS A110058. With the classical leading equivalent
\[
L_n=\frac{e^{1/4}4^{n^2}}{(4\pi)^{n-1/2}n^{n-1}},
\]
we prove that, for every fixed $K\geq0$, the logarithm $\log(a_n/L_n)$ has a rational inverse-power expansion through $n^{-K}$ with error $O_K(n^{-K-1})$. A finite connected-Wick algorithm determines every coefficient; the first two are $-3/2$ and $223/32$. The proof combines Canfield--McKay's absolute localization theorem for integer tables with Isaev's complex cumulant theorem. We give the gauge normalization, product cutoff, all-order variation bounds, and connected-diagram power counting needed for this application. An independent formal Gaussian factorization verifies all eighteen cumulants entering the first two corrections. A compact-positive-density extension has coefficient polynomials in $1/[\lambda(1+\lambda)]$ and proves the inequalities $0<\Delta<2$ of the Canfield--McKay correction-factor conjecture in one regime only: square tables whose density $s/n$ lies in a fixed compact subset of $(0,\infty)$, for all $n$ beyond a threshold that is not made explicit. The conjecture itself, stated for all sizes, shapes and densities, is not settled. The expansion also gives a two-ceiling enclosure for the inverse counting function, with real-index uncertainty $O(t^{-4})$ at two corrections. Constants in the asymptotic remainder are not made numerically explicit. The accompanying exact computations verify identities and small counts; they are not finite-size error certificates.
\end{abstract}

\begin{writenote}[Status and trust boundary, added 6 October 2026, batch~108]
This is bundle Report~236 of an external AI-assisted research session, filed
in ProveIt as the single-source research report
\file{a110058-square-contingency-tables} (provenance, the sources as the
write read them, relation to the repository, reading conventions and the
collected non-claims in Section~\ref{sqt:sec:provenance}). It is unrefereed
and not formalized; its author line and PDF author field read ``Report 236''.
\emph{Narrowed wording.} The abstract's sentence on the correction-factor
conjecture read, as delivered, ``A compact-positive-density extension has
coefficient polynomials in $1/[\lambda(1+\lambda)]$ and proves an eventual
square case of the Canfield--McKay correction-factor conjecture.'' Canfield
and McKay conjecture $0<\Delta<2$ for \emph{every} admissible 4-tuple;
Corollary~\ref{sqt:cor:delta} proves it for the square tuples
$(n,s;n,s)$ with $s/n$ in a fixed compact $I\subset(0,\infty)$ and
$n\ge n_0(I)$, with $n_0(I)$ not explicit, so it settles one regime, not the
conjecture, and verifies it at no explicitly named tuple
(Remark~\ref{sqt:rem:cmregime}, which quotes both statements). The
sentence was narrowed accordingly, as were the title of
Section~\ref{sqt:sec:conjecture} and one sentence after the corollary (first
wordings kept there). \emph{What rests on what.} The analytic proof imports
two theorems, cited and not reproved: Canfield--McKay's absolute
localization (Theorem~3 of arXiv:math/0703600v2) and Isaev's complex
cumulant theorem (Theorem~1.2 of arXiv:2508.16952v2); the write read both
statements and checks the specializations in
Section~\ref{sqt:sec:provenance}. For the density extension
(Theorem~\ref{sqt:thm:density}) the source asserts that Canfield--McKay's
localization holds \emph{uniformly} for $\lambda$ in a compact interval and
supports this by its own reading of their proof; that is a trust boundary of
the source. The write adds Lemma~\ref{sqt:lem:sequential}, which derives the
uniform superpolynomial smallness that Theorem~\ref{sqt:thm:density} uses
from the \emph{statement} of their Theorem~3 alone, and reports what the
write found in their proof (Remark~\ref{sqt:rem:uniformity},
Question~\ref{sqt:q:uniform}). Coefficient values are checked by exact finite
computations in the package; no proof uses a floating-point computation. The
write also adds Remark~\ref{sqt:rem:oeis} (the OEIS entry, whose only
asymptotic formula is Pak's log-scale one),
Remark~\ref{sqt:rem:transseries} (Section~\ref{sqt:sec:inverse} against the
transseries volume) and Section~\ref{sqt:sec:further} (open questions).
\end{writenote}
\setcounter{tocdepth}{1}
\begingroup
\small\setlength{\parskip}{0pt}
\makeatletter
\renewcommand\l@section[2]{\@dottedtocline{1}{0em}{2em}{#1}{#2}}
\makeatother
\tableofcontents
\endgroup
\clearpage
\section{The result and its scope}\label{sqt:sec:results}
For $n\geq1$, define
\[
a_n=\#\left\{(m_{ij})\in\mathbb Z_{\geq0}^{n\times n}:
 \sum_jm_{ij}=n\text{ for all }i,\quad
 \sum_im_{ij}=n\text{ for all }j\right\}.
\]
Rows and columns are labelled. Thus matrices related by a permutation of rows or columns are generally distinct. The first values are $1,3,55,10147,22069251$ for $n=1,\ldots,5$; with $a_0=1$, this is OEIS A110058 \cite{oeis}. Throughout, all logarithms are natural, and asymptotic assertions concern integer $n\to\infty$ unless a real variable is specified. Set
\begin{equation}\label{sqt:eq:leading}
G_n=\frac{4^{n^2}}{(4\pi)^{n-1/2}n^{n-1}},\qquad L_n=e^{1/4}G_n.
\end{equation}
The leading equivalent $a_n\sim L_n$ is already contained in Canfield--McKay \cite{CM}; it is also covered by the smooth-margin theory of Barvinok--Hartigan \cite{BH}. The purpose here is to justify every fixed algebraic order in this square density-one specialization, evaluate two corrections, and control inversion.

\begin{theorem}\label{sqt:thm:main}
There are rational numbers $C_j$, $j\geq1$, such that for every fixed integer $K\geq0$,
\begin{equation}\label{sqt:eq:main}
\log\frac{a_n}{L_n}=\sum_{j=1}^{K}\frac{C_j}{n^j}+O_K(n^{-K-1}).
\end{equation}
Each $C_j$ is given by a finite rational combinatorial algorithm. In particular,
\begin{equation}\label{sqt:eq:two}
C_1=-\frac32,\qquad C_2=\frac{223}{32}.
\end{equation}
Consequently,
\begin{align}
\log\frac{a_n}{L_n}&=-\frac{3}{2n}+\frac{223}{32n^2}+O(n^{-3}),\label{sqt:eq:logtwo}\\
\frac{a_n}{L_n}&=1-\frac{3}{2n}+\frac{259}{32n^2}+O(n^{-3}).\label{sqt:eq:relativetwo}
\end{align}
\end{theorem}

The effective assertion is about individual coefficients. It does not assert explicit values of the implicit remainder constants or starting indices. The theorem is a Poincar\'e expansion at every fixed order, with $K$ fixed before $n\to\infty$. It does not assert convergence of the infinite series, uniformity in growing $K$, optimal truncation, or a description of exponentially small contributions.

The argument has three distinct parts. First, an exact Fourier integral is localized by the integer-table theorem of Canfield--McKay. A Gaussian gauge penalty and a linear whitening map then convert the localized problem into an expectation on a product of truncated Gaussian intervals. Second, Isaev's theorem for bounded \emph{complex} functions \cite{Isaev} gives a sufficiently accurate finite cumulant expansion. Third, connected Wick contractions provide exact Laurent polynomials and a degree bound that leaves only finitely many degree multisets at any required order. The complex cumulant method is imported; the checks needed to apply it to these geometric tables are supplied in full.

For later use, define the inverse counting function
\begin{equation}\label{sqt:eq:inversecount}
\nu(Y)=\min\{n\geq1:a_n\geq Y\},\qquad Y>1,
\end{equation}
and, for fixed $K$, the real approximation
\begin{equation}\label{sqt:eq:FK}
F_K(x)=x^2\log4-(x-\tfrac12)\log(4\pi)-(x-1)\log x
       +\frac14+\sum_{j=1}^{K}C_jx^{-j}.
\end{equation}
There is a unique sufficiently large solution $t_K=t_K(Y)$ of $F_K(t_K)=\log Y$. We prove in Section~\ref{sqt:sec:inverse} that some $D_K>0$ gives, for all sufficiently large $Y$,
\begin{equation}\label{sqt:eq:inverseintro}
\ceil{t_K-D_Kt_K^{-K-2}}\leq\nu(Y)\leq
\ceil{t_K+D_Kt_K^{-K-2}}.
\end{equation}
The two ceilings matter near an integer threshold. An unqualified rule $\nu(Y)=\lceil t_K\rceil$ is not a consequence of the theorem.

\begin{remark}[{[write]} The OEIS entry, 6~October 2026]\label{sqt:rem:oeis}
The live entry A110058 \cite{oeis} (revision~\#30 of 30~May 2026, read
6~October 2026) is named ``Number of nonnegative integer matrices of order n
for which all row and column sums equal n.'' Its data run through $a_{13}$,
its one link is to the Combinatorica version of \cite{CM}, and its only
formula, the only asymptotic statement in the entry, reads verbatim
\begin{center}\small
\texttt{log a(n) = 2(log 2)*n\^{}2 - n*(log n) - n*(log 4*Pi) + (log n) + O(1). - \_Igor Pak\_, May 15 2019}
\end{center}
(with \texttt{log 4*Pi} read as $\log(4\pi)$). Since
$\log L_n=n^2\log4-(n-\tfrac12)\log(4\pi)-(n-1)\log n+\tfrac14$, Pak's main
terms are exactly $\log L_n-\tfrac12\log(4\pi)-\tfrac14$. So the formula is
consistent with $a_n\sim L_n$, which is Canfield--McKay's, and the latter
identifies Pak's bounded term as $\tfrac12\log(4\pi)+\tfrac14+o(1)
=1.51551\ldots+o(1)$; Theorem~\ref{sqt:thm:main} expands the $o(1)$ to every
fixed order. Nothing in the entry is corrected, and nothing was submitted to
the OEIS. The ``b-file'' cited by the source is synthesized by the OEIS from
the entry's data (its header reads ``A110058 (b-file synthesized from
sequence entry)'', $0\le n\le13$); the counts at $n=6,8,10,13$ in
Table~\ref{sqt:tab:diagnostics} are taken from it.
\end{remark}

\subsection{Provenance, sources and reading conventions}\label{sqt:sec:provenance}
\begin{writenote}[Added 6 October 2026, batch~108]
\emph{Provenance.} This report is Report~236 of the session bundle that
ProveIt's \file{docs/incoming} intake processed as batch~108: the archive
\file{Report236.zip} (669{,}911 bytes, SHA-256
\file{1b73cd6bd55d...5b132059}, 35 files in a wrapper directory
\file{Report236/}), arrival commit \file{60f54ea06}, placed in
\file{602e5bd0f}. The text is the delivered \file{article.tex} with its
twelve files \file{sections/*.tex} (980 lines in all, dated 5~October 2026),
shipped under the same names. The delivered 24-page PDF and the checksum
manifest \file{MANIFEST.sha256} (34 entries, verified at placement) are not
shipped, and the delivered README is replaced by the report README; all
three are retrievable from the arrival commit. The package records no
ProveIt commit and cites nothing in the repository. The phrase ``private
review'' occurs in it only in sentences saying that such material is not
included (Section~\ref{sqt:sec:computation}, \file{SOURCES.md} and the
delivered README); the PDF author field is ``Report 236''. It is a
single-source report, so no merge choice was made. The write prefixed the
97 delivered labels with \file{sqt:} and updated the 87 references to them;
it added this subsection, the notes marked [write],
Remarks~\ref{sqt:rem:oeis}, \ref{sqt:rem:uniformity},
\ref{sqt:rem:cmregime} and~\ref{sqt:rem:transseries},
Lemma~\ref{sqt:lem:sequential} and Section~\ref{sqt:sec:further}, and it
narrowed three sentences on the correction-factor conjecture (the abstract,
the title of Section~\ref{sqt:sec:conjecture} and the sentence after
Corollary~\ref{sqt:cor:delta}; first wordings kept in the notes). No
theorem, proof or number of the source was changed: the added statements are
the last of their sections, the added section comes after the last numbered
one, and the added displays are unnumbered, so every theorem, section,
table and equation keeps its delivered number.

\emph{The sources, as the write read them} (6~October 2026).
\begin{itemize}\raggedright
\item OEIS A110058, revision~\#30, and its synthesized b-file
 (Remark~\ref{sqt:rem:oeis}).
\item Canfield--McKay \cite{CM}, arXiv:math/0703600v2: pages 1--6 and
 15--20 (Theorem~1 with condition~(1.1), Conjecture~1 with~(1.4) and their
 record of proved cases, the convention on asymptotic constants, the
 integral~(2.2)--(2.3), the region~(3.1), Theorem~3, Lemmas~3 and~4 and
 the proof of Theorem~3). Checked against the source: at $m=n$, $s=t=\lambda
 n$ their integrand $F$ of~(2.3) is $\mathcal F_\lambda$ of
 Section~\ref{sqt:sec:density}, since
 $e^{-\ii s\sum_j\theta_j-\ii t\sum_k\phi_k}=\prod_{j,k}e^{-\ii\lambda(\theta_j+\phi_k)}$;
 their $I_0$ of~(4.1) is $I_{0,\lambda}$ of~\eqref{sqt:eq:densityminor}
 (and $I_0$ of~\eqref{sqt:eq:minor} at $\lambda=1$); their region~(3.1) is
 the region described in Sections~\ref{sqt:sec:imports}
 and~\ref{sqt:sec:density}, with their $\varepsilon$ equal to the source's
 $\eta$; the left side of~(1.1) is $3$ at $\lambda=1$ and
 $\tfrac23(4+1/u)$ in general; and their leading formula~(1.3) at $m=n$ is
 $\mathcal G(n,\lambda)e^{B_0(u)}$, so $B_0$ and $\mathcal G$ (and $L_n$ at
 $\lambda=1$) are theirs. Their Theorem~3 gives $O(e^{-n^\varepsilon})I_0$;
 Proposition~\ref{sqt:prop:minor} states the weaker $O(e^{-cn^\eta})I_0$.
 Not read: their Sections~3 and~5 (the central evaluation and the proof of
 Theorem~1), which the source does not use.
\item Isaev \cite{Isaev}, arXiv:2508.16952v2 (HTML version): the statement
 of Theorem~1.2 and the definitions of $\Delta_V$, $\overline\Delta_V$ and
 $S_k$. Proposition~\ref{sqt:prop:complex} is that theorem with $n=N$ and
 $S_k$ replaced by the larger $D_k$, since
 $\overline\Delta_V(f,X)\le\Delta_V(f)$.
\item The transseries volume
 \path{Analysis/Transseries/docs/series-and-transseries/Transseries_And_Inversion/transseries_and_inversion.tex}:
 \file{p0:def:core}, \file{p0:thm:lambert-core},
 \file{p0:prop:factorial-core}, \file{p0:thm:core-reversion},
 \file{p0:def:three-inverses}, \file{p0:thm:staircase} and
 \file{plt:thm:lw-template}, read for Remark~\ref{sqt:rem:transseries}; and,
 at the independent check of 6~October 2026,
 \file{plt:def:lw-monomial-log-datum}.
\item Not read by the write: Barvinok--Hartigan \cite{BH} (beyond its arXiv
 title page), Greenhill--McKay \cite{GM}, Isaev--McKay \cite{IM,IM2018},
 Isaev--McKay--Zhang \cite{IMZ}, Zipunnikov--Booth--Yoshida \cite{ZBY}, Aw
 \cite{Aw} and Isaev--Makai--McKay \cite{IMM}. What is said about them
 (Barvinok--Hartigan's Theorem~1.3 and the comparison~\eqref{sqt:eq:BHcompare},
 the hypotheses of Isaev--McKay, Greenhill--McKay's Corollaries~4.1
 and~4.2) is the source's. Canfield--McKay's own record (their case~(b),
 quoted in Remark~\ref{sqt:rem:cmregime}) independently states that the
 conjecture holds eventually in Greenhill--McKay's sparse range.
\end{itemize}

\emph{What was checked at intake.} The batch-108 dossier read the
manuscript, checked the conjecture and condition~(1.1) against the
Canfield--McKay text, and reran the package on a copy (Windows, Python
3.14.4, SymPy 1.14.0): \file{build.py --verify-only} verified 34 files;
\file{connected_wick.py}, \file{verify_coefficients.py} and
\file{count_tables.py} reproduced the delivered receipts after line-ending
conversion; \file{diagnostics.py} reproduced
Table~\ref{sqt:tab:diagnostics}; \file{guard_tests.py} stops on Windows
at its own path guard; the PDF and ZIP rebuilds were not run. It
computed $\Delta(n,n;n,n)$ from exact counts and found $1.070,\ldots,0.565$
for $n=4,5,6,8,10,13$, and Richardson extrapolation of $D_1,D_2$ over
$6\le n\le13$ gave about $-1.47$ and $6.5$ against $-3/2$ and $223/32$. The
write rechecked by hand the coefficient identity
$P_2(1/u)-g_2(u)-\tfrac56=\tfrac{67}{12}+\tfrac5{6u}$, the expansion of
$(n-1)\log(1+1/n)$, the two normalization identities of
Section~\ref{sqt:sec:density}, the value $3$ of condition~(1.1) and the
identifications listed above; it computed $\Delta(n,n;n,n)$ for all
$1\le n\le13$ (Remark~\ref{sqt:rem:cmregime}) and checked in SymPy that the
four reversion coefficients of~\eqref{sqt:eq:reversion} solve the
fixed-point equation of Remark~\ref{sqt:rem:transseries}. Finite
computations are observations, not certificates.

\emph{Relation to the repository.} Before batch~108 no file of the
repository mentioned A110058 or Canfield and McKay's conjecture. No
statement of this report is formalized in Lean or Rocq, and its place in the
collection confers no formal status. The neighbouring reports are listed in
the report README; \file{a380592-tied-football-seasons} applies the same
Isaev theorem to another integral and receives a reciprocal note.
\end{writenote}
\begin{writenote}[What is not claimed, collected from the source]
No numerical value of a remainder constant or a starting index is asserted:
Theorem~\ref{sqt:thm:main} is a Poincar\'e expansion at each fixed order $K$,
with no convergence of the series, no uniformity in growing $K$, no optimal
truncation and no exponentially small terms; no nonperturbative sectors,
resurgence or Stokes behaviour are established. The leading equivalent is
Canfield--McKay's and is also within Barvinok--Hartigan's theory; Isaev's
complex cumulant theorem and the connected-Wick toolkit are imported, and
no new general saddle-point or cumulant method is claimed. The literature
review was bounded: absence of the coefficient pair from the sources
inspected ``is not a worldwide priority certificate''; Greenhill--McKay's
sparse result is the antecedent of the first coefficient, and
equation~\eqref{sqt:eq:BHcompare} is a specialization of Barvinok--Hartigan's
formula, not a correction. The inverse constants are existential:
$\lceil t_K\rceil$ is not a valid rule, computing a root does not compute
$B_K$, and the Newton statement is about exact-real iteration. The density
extension is not uniform as the density tends to $0$ or~$\infty$; at a fixed
irrational density there are no exact margins. The corollary on the
correction factor has no explicit starting index and does not settle finite
sizes, nonsquare tables or degenerating densities. No monotonicity of the
diagnostics or sign of omitted terms is claimed. The package's computations
verify finite identities and small counts in a fixed validated range (cost at
most three); a receipt certifies no remainder, localization, finite-size
inequality, inverse ceiling or priority; byte-identical PDF reproduction is
toolchain-specific; the manifest is ``not a digital signature''; and the
reproducibility controls are ``not a sandbox''. The write adds: the
uniformity in density of the imported localization is established in this
text only in the weaker form of Lemma~\ref{sqt:lem:sequential}
(Question~\ref{sqt:q:uniform}).
\end{writenote}
\medskip\noindent\begin{minipage}{\textwidth}
\emph{Reading conventions} (letters the manuscript reuses, with the tempting
false readings; no symbol of the source was renamed; symbols of the
transseries volume carry the subscript ``vol'' in
Remark~\ref{sqt:rem:transseries}).\par
\smallskip
\begingroup\small\renewcommand{\arraystretch}{1.1}
\noindent\begin{tabular}{@{}>{\raggedright\arraybackslash}p{0.85in}>{\raggedright\arraybackslash}p{3.15in}>{\raggedright\arraybackslash}p{2.15in}@{}}
\toprule
Symbol & Meaning here (where) & Not to be read as\\
\midrule
$\Delta$ & Isaev's mixed difference $\Delta_V(f)$ (Section~\ref{sqt:sec:imports}); Canfield--McKay's correction $\Delta(m,s;n,t)$ (Section~\ref{sqt:sec:conjecture}) & each other\\
$D_k$, $D_K$, $D_1$, $D_2$ & Isaev's sums $D_k(f)$; the constant of Theorem~\ref{sqt:thm:inverse}; the diagnostics of Section~\ref{sqt:sec:computation} & each other\\
$G_n$, $G_{\rm Good}$, $\mathcal G$ & Gaussian prefactor~\eqref{sqt:eq:leading}; Canfield--McKay's (Good's) binomial estimate; density prefactor~\eqref{sqt:eq:densityG} & each other\\
$A$, $A_\lambda$ & augmented matrix~\eqref{sqt:eq:augmented}; the scalar $u/2$, Canfield--McKay's $A$ & each other\\
$T$, $T(n,s)$ & whitening map~\eqref{sqt:eq:whitening}; the table count of Section~\ref{sqt:sec:density} & each other\\
$h$ & local interaction~\eqref{sqt:eq:h}; the gauge coordinate $x_{2n}$ (Section~\ref{sqt:sec:gauge}) & $h_{\rm vol}$ of \file{p0:thm:core-reversion}\\
$m$, $L$ & cumulant order and Taylor cutoff~\eqref{sqt:eq:cutofforders} & Canfield--McKay's row count $m$ ($=n$ here); $L_n$ of~\eqref{sqt:eq:leading}\\
$\varepsilon$, $\eta$ & cube exponent; localization exponent $\eta=\varepsilon/2$ & Canfield--McKay's $\varepsilon$, which is the source's $\eta$\\
$a$, $b$ & in Section~\ref{sqt:sec:imports}: the constants of condition~(1.1); in Section~\ref{sqt:sec:inverse}: $\log4$, $\log(4\pi)$ & each other; $a_n$\\
$r$ & number of cumulant arguments; $\sqrt{\log Y/a}$ (Section~\ref{sqt:sec:inverse}); exceptional count (Section~\ref{sqt:sec:density}) & each other\\
$u$ & $\lambda(1+\lambda)$; the gauge parameter in the proof of Lemma~\ref{sqt:lem:transfer} & $u_{\rm vol}$ of \file{p0:thm:core-reversion}\\
$d,e,f,g$ & reversion coefficients~\eqref{sqt:eq:reversiond}--\eqref{sqt:eq:reversiong} & degrees $d_a$; $f=R\circ T$\\
$P_j$, $\Part$ & density polynomials; set partitions & each other; $P_a(A)$ power sums\\
$\nu(Y)$ & $\min\{n\ge1:a_n\ge Y\}$ & it \emph{is} the staircase $N_\ast$ of \file{p0:def:three-inverses} (Remark~\ref{sqt:rem:transseries})\\
$\kappa$ & cumulants $\kappa_r$; geometric cumulants $\kappa_d(\lambda)$ & the $\kappa$ of \file{p0:prop:factorial-core}\\
\bottomrule
\end{tabular}
\endgroup
\end{minipage}\par

\section{The two imported analytic results}\label{sqt:sec:imports}
This section states exactly the consequences used from the primary sources. Their proofs are not reproduced here. The remaining arguments, including the reduction to their hypotheses, are given below.

\subsection{Absolute localization for geometric tables}
For $\theta,\phi\in(\R/2\pi\mathbb Z)^n$, put
\begin{equation}\label{sqt:eq:integrand}
\psi(z)=\frac{e^{-\ii z}}{2-e^{\ii z}},\qquad
\mathcal F(\theta,\phi)=\prod_{i,j=1}^n\psi(\theta_i+\phi_j).
\end{equation}
Fix a sufficiently small positive $\eta$. In Canfield--McKay's region $\mathcal R_\eta$, each angular vector lies in a semicircle, the centered row and column deviations have absolute value at most $\tfrac12n^{-1/2+\eta}$, and the circular sum $\mu$ of the two angular means obeys $|\mu|\leq\tfrac12n^{-1+4\eta}$. Means and deviations are evaluated in the corresponding small circular chart.

\begin{proposition}[Canfield--McKay, specialized]\label{sqt:prop:minor}
For all sufficiently large $n$,
\begin{equation}\label{sqt:eq:minor}
\int_{\mathcal R_\eta^c}|\mathcal F(\theta,\phi)|\,\dd\theta\,\dd\phi
 =O(e^{-c n^\eta})I_0,
\qquad I_0=\pi^{n-1/2}n^{1-n}e^{-1/4},
\end{equation}
where $c>0$ can be taken fixed.
\end{proposition}
This is Theorem~3 with the region in equation~(3.1) of the preprint \cite{CM}. At $m=n$, line sums $s=t=n$, and density $\lambda=1$, the left side of its condition~(1.1) is
\[
\frac{(1+2\lambda)^2}{4\lambda(1+\lambda)}
\left(1+\frac{5m}{6n}+\frac{5n}{6m}\right)=3.
\]
Thus the condition $3\leq a\log n$, with fixed $a,b>0$ and $a+b<1/2$, holds eventually. The polynomial upper bound on $\lambda$ in Theorem~3 is automatic. These are asymptotic conditions; no numerical starting index is inferred. We use the absolute minor-arc bound, not the coarser central Taylor remainder in the leading enumeration formula.

\subsection{A complex cumulant tail estimate}
Let $X=(X_1,\ldots,X_N)$ have independent coordinates taking values in a product domain $\Omega=\prod_{j=1}^N\Omega_j$. If $x,y\in\Omega$ and $V\subseteq[N]$, write $x\mathbin{\triangleright_V}y$ for the vector obtained by replacing the coordinates of $x$ in $V$ by those of $y$. For bounded measurable $f:\Omega\to\mathbb C$, define
\begin{align}
\Delta_V(f)&=\sup_{x,y\in\Omega}\left|
 \sum_{W\subseteq V}(-1)^{|W|}f(x\mathbin{\triangleright_W}y)\right|,\\
D_k(f)&=\max_{j\in[N]}\sum_{\substack{|V|=k\\j\in V}}\Delta_V(f).
\end{align}
Isaev's distribution-dependent quantities use
\[
\Delta_V(f,X)=\sup_{y\in\Omega}\left|\E\sum_{W\subseteq V}
 (-1)^{|W|}f(X\mathbin{\triangleright_W}y)\right|,
\]
and the corresponding sums $S_k(f,X)$. In particular, $S_k(f,X)\leq D_k(f)$. Cumulants of complex variables use ordinary products without complex conjugation.

\begin{proposition}[Isaev, sufficient form]\label{sqt:prop:complex}
Let $m$ be a positive integer. Suppose $f$ is bounded and measurable and, for every $1\leq t\leq m$,
\begin{equation}\label{sqt:eq:isaevhyp}
\sum_{k=t}^N e^{3\alpha k/2}2^t\binom{k-1}{t-1}D_k(f)
 \leq\alpha\leq\frac1{100}.
\end{equation}
Then there exists a complex $\delta$ such that
\begin{equation}\label{sqt:eq:isaev}
\E e^{f(X)}=(1+\delta)^N
 \exp\left(\sum_{r=1}^m\frac{\kappa_r(f(X))}{r!}\right),
\qquad |\delta|\leq e^{(100\alpha)^{m+1}}-1.
\end{equation}
\end{proposition}
This follows directly from Theorem~1.2 of \cite{Isaev} by $S_k\leq D_k$. The factor $(1+\delta)^N$ is important: at fixed $m$, if $N\alpha^{m+1}=o(1)$, the total multiplicative error is $1+O_m(N\alpha^{m+1})$. Bounds on mixed differences through order $m$ alone would not suffice; the sums in \eqref{sqt:eq:isaevhyp} continue to $N$.

The leading equivalent can alternatively be obtained from Barvinok--Hartigan's smooth-margin theorem \cite[Theorem~1.3]{BH}. Their absolute localization theorem is also compatible with the argument, but is not needed here. The binary graph-factor enumeration theorem of Isaev--McKay \cite{IM} is not used: its signless-Laplacian assumptions exclude a bipartite base graph. Nor does an enumeration theorem for Eulerian orientations \cite{IMZ} apply to the geometric entries. The distinction between Isaev's complex theorem and its earlier real-valued predecessor is essential because our cubic term is imaginary.

\section{The exact integral and a product cutoff}\label{sec:gauge}
\subsection{Fourier inversion and the null direction}
Let the $n^2$ entries $X_{ij}$ be independent with
$\PP(X_{ij}=k)=2^{-k-1}$ for $k\geq0$. Every matrix with the prescribed margins has total sum $n^2$ and probability $4^{-n^2}$. The characteristic function of $X_{ij}-1$ is $\psi$ in \eqref{eq:integrand}. Fourier inversion of all $2n$ margin constraints, including their one redundancy, gives the exact identity
\begin{equation}\label{eq:fourier}
a_n=4^{n^2}(2\pi)^{-2n}\mathcal J_n,\qquad
\mathcal J_n=\int_{[-\pi,\pi]^{2n}}\mathcal F(\theta,\phi)\,\dd\theta\,\dd\phi.
\end{equation}
The invariance under $(\theta,\phi)\mapsto(\theta+t\1,\phi-t\1)$ is the only continuous gauge direction.

Write $x=(\theta,\phi)$ and $v=(\1_n,-\1_n)$. To avoid confusing the integral with an all-ones matrix, denote the latter by $\mathbf J$. The quadratic term is
\[
q(x)=\sum_{i,j}(x_i+x_{n+j})^2=x^{\mathsf T}Qx,
\qquad Q=\begin{pmatrix}nI&\mathbf J_n\\\mathbf J_n&nI\end{pmatrix}.
\]
Since $Q+vv^{\mathsf T}/2=nI+\mathbf J_{2n}/2$, define
\begin{equation}\label{eq:augmented}
q_+(x)=q(x)+\frac{(v\cdot x)^2}{2}=nx^{\mathsf T}Ax,
\qquad A=I+\frac{\mathbf J_{2n}}{2n}.
\end{equation}
The eigenvalues of $A$ are $2$ on the all-ones direction and $1$ on its orthogonal complement. Hence
\begin{equation}\label{eq:whitening}
T=A^{-1/2}=I+\gamma\frac{\mathbf J_{2n}}{2n},
\qquad\gamma=\frac1{\sqrt2}-1,\qquad\det T=\frac1{\sqrt2}.
\end{equation}
Both $T$ and $T^{-1}$ have uniformly bounded $1$- and infinity-operator norms, and $Tv=T^{-1}v=v$. The augmented Gaussian integral is
\begin{equation}\label{eq:Zn}
Z_n=\int_{\R^{2n}}e^{-q_+(x)}\,\dd x=\frac1{\sqrt2}\left(\frac\pi n\right)^n.
\end{equation}

\subsection{The exact gauge factor}
Use real coordinates
\[
s_i=x_i+x_{2n}\quad(1\leq i\leq n),\qquad
t_j=x_{n+j}-x_{2n}\quad(1\leq j<n),\qquad h=x_{2n}.
\]
Their inverse is $x_i=s_i-h$, $x_{n+j}=t_j+h$, $x_{2n}=h$, and its absolute Jacobian is $1$. Pair sums, and therefore $\mathcal F$, are independent of $h$. Moreover,
\[
v\cdot x=\sum_i s_i-\sum_{j<n}t_j-2nh,
\qquad
\int_{\R}e^{-(v\cdot x)^2/2}\,\dd h=\frac{\sqrt{2\pi}}{2n}.
\]
On the torus, setting $\phi_n=0$ instead leaves a gauge orbit factor $2\pi$. Thus the conversion from an augmented real integral on a full lift to the corresponding torus integral multiplies by
\begin{equation}\label{eq:gaugefactor}
\frac{2\pi}{\sqrt{2\pi}/(2n)}=2n\sqrt{2\pi}.
\end{equation}
This also verifies the eventual prefactor directly:
\begin{equation}\label{eq:prefactor}
\frac{2n\sqrt{2\pi}\,Z_n}{(2\pi)^{2n}}
 =(4\pi)^{-n+1/2}n^{1-n}.
\end{equation}

\subsection{Transfer to a whitened cube}
Choose $0<\eps\leq1/100$ small enough that Proposition~\ref{prop:minor} applies with $\eta=\eps/2$. Put
\begin{equation}\label{eq:cubes}
\rho=n^{-1/2+\eps},\qquad
\mathcal C_n=[-\rho,\rho]^{2n},\qquad\mathcal B_n=T\mathcal C_n.
\end{equation}
\begin{lemma}\label{lem:transfer}
For some fixed $c>0$,
\begin{equation}\label{eq:transfer}
\mathcal J_n=2n\sqrt{2\pi}\left\{
\int_{\mathcal B_n}\mathcal F(x)e^{-(v\cdot x)^2/2}\,\dd x
 +O(e^{-c n^\eta})Z_n\right\}.
\end{equation}
\end{lemma}
\begin{proof}
Each orbit in $\mathcal R_\eta$ has a small representative $x_H\in v^\perp$ given, in its local real chart, by
\[
x_{H,i}=\theta_i-\bar\theta+\mu/2,\qquad
x_{H,n+j}=\phi_j-\bar\phi+\mu/2.
\]
Consequently,
\begin{equation}\label{eq:smallrep}
\|x_H\|_\infty\leq\tfrac12n^{-1/2+\eta}+\tfrac14n^{-1+4\eta}=o(\rho),
\qquad \|T^{-1}x_H\|_\infty=o(\rho).
\end{equation}
The associated $s,t$ coordinates lie in $(-\pi,\pi)$ for large $n$. These small representatives are injective in the quotient torus: if two differ by a projected $2\pi$ integer vector, their row-plus-column differences are multiples of $2\pi$ of magnitude $o(1)$, hence zero. Their difference is then a gauge vector, and membership in $v^\perp$ makes it zero.

Let $D$ be the resulting small $s,t$ region and $\widetilde D$ its full real lift, with $h$ unrestricted. The exact gauge calculation gives
\begin{equation}\label{eq:liftidentity}
\int_{\mathcal R_\eta}\mathcal F
 =2n\sqrt{2\pi}\int_{\widetilde D}\mathcal F(x)e^{-(v\cdot x)^2/2}\,\dd x.
\end{equation}
We compare $\widetilde D$ and $\mathcal B_n$ in the two directions separately.

On $\mathcal B_n\setminus\widetilde D$, all coordinates are $O(\rho)$, so the quotient projection is still in an injective small chart and outside $D$. The integral of the nonnegative gauge weight over any subset of a fiber is at most the full fiber integral $\sqrt{2\pi}/(2n)$. Thus the absolute integral on this difference is at most $(2n\sqrt{2\pi})^{-1}$ times the torus minor-arc integral. By \eqref{eq:minor}, this is $O(e^{-c n^\eta})Z_n$, since $I_0/(2n\sqrt{2\pi}Z_n)=e^{-1/4}/(2\pi)$ is constant.

On $\widetilde D\setminus\mathcal B_n$, write $x=x_H+uv$. By \eqref{eq:smallrep} and $T^{-1}v=v$, exclusion from $\mathcal B_n$ forces $|u|>\rho/2$ eventually. Since $v\cdot x=2nu$, the omitted gauge fiber has relative Gaussian mass at most $O(e^{-c n^2\rho^2})$. Here $n^2\rho^2=n^{1+2\eps}$. Even using only $|\mathcal F|\leq1$, quotient volume at most $(2\pi)^{2n-1}$, and the normalization \eqref{eq:Zn}, the loss is bounded by
\[
Z_n\exp\{-c n^{1+2\eps}+O(n\log n)\}
 =O(e^{-c' n^\eta})Z_n.
\]
Finally, the torus outside $\mathcal R_\eta$ has the absolute bound \eqref{eq:minor}. Combining this with \eqref{eq:liftidentity} proves \eqref{eq:transfer}.
\end{proof}

The asymmetric choice $\eta=\eps/2$ leaves the main quotient region strictly inside the whitened cube. It avoids a comparison that would multiply a Gaussian tail by a crude supremum of the positive quartic exponential.

\section{The analytic cumulant reduction}\label{sqt:sec:analytic}
\subsection{Keep the exact remainder on a bounded domain}
Take the analytic logarithm equal to zero at the origin and write
\begin{equation}\label{sqt:eq:h}
h(z)=-\ii z-\log(2-e^{\ii z})+z^2=\sum_{d\geq3}h_dz^d,
\qquad R(x)=\sum_{i,j}h(x_i+x_{n+j}).
\end{equation}
The closest zero of $2-e^{\ii z}$ is $-\ii\log2$, so $h$ is analytic for $|z|<\log2$. Cauchy estimates on a smaller fixed disk give constants independent of $n$ such that, on $|z|\leq C\rho$,
\begin{equation}\label{sqt:eq:derivativebounds}
h'(z)=O(\rho^2),\quad h''(z)=O(\rho),\quad
|h^{(k)}(z)|\leq k!C_0^k\quad(k\geq3).
\end{equation}
The first coefficients are
\begin{equation}\label{sqt:eq:hsix}
h_3=-\ii,\quad h_4=\frac{13}{12},\quad h_5=\frac{5\ii}{4},\quad
h_6=-\frac{541}{360},\quad h_7=-\frac{223\ii}{120},\quad
h_8=\frac{47293}{20160}.
\end{equation}
More generally, if $\genfrac{\{}{\}}{0pt}{}{a}{b}$ is a Stirling number of the second kind,
\begin{equation}\label{sqt:eq:stirling}
h_d=\frac{2\ii^d}{d!}\sum_{j=1}^{d-1}\genfrac{\{}{\}}{0pt}{}{d-1}{j}j!,\qquad d\geq3.
\end{equation}
Indeed the $d$th geometric cumulant is $\sum_{k\geq1}k^{d-1}2^{-k}$, and the falling-factorial expansion of $k^{d-1}$ evaluates that sum as twice the finite Stirling sum. In particular all $h_d/\ii^d$ are rational.

Let $X$ have $2n$ independent $N(0,1/(2n))$ coordinates, and let $\widehat X$ be $X$ conditioned on $\mathcal C_n$. The conditioning event is a product of intervals, so the coordinates remain independent. With $f=R\circ T$,
\begin{equation}\label{sqt:eq:boundedexpectation}
\int_{\mathcal B_n}\mathcal F(x)e^{-(v\cdot x)^2/2}\,\dd x
 =Z_n\PP(X\in\mathcal C_n)\E e^{f(\widehat X)},
\end{equation}
where
\begin{equation}\label{sqt:eq:tailprob}
\PP(X\notin\mathcal C_n)\leq O(n)e^{-n^{2\eps}}
 =O(e^{-c n^{2\eps}}).
\end{equation}
All exponential expectations in this proof use this bounded product domain. Full Gaussians will appear only in fixed polynomial moments and cumulants.

\subsection{Control of every mixed-difference order}
For a cell $e=(i,j)$ write its whitened pair sum as $a_e\cdot y=(Ty)_i+(Ty)_{n+j}$. Its coefficient vector has two entries $1+\gamma/n$ and $2n-2$ entries $\gamma/n$. With $\beta=-\gamma=1-1/\sqrt2$,
\begin{equation}\label{sqt:eq:rowsumbounds}
\sum_{q=1}^{2n}|a_{e,q}|\leq2+2\beta,\qquad
\sum_e|a_{e,q}|\leq(1+\beta)n.
\end{equation}
For a coordinate set $V$ of size $k$, repeated integration of the mixed derivative along the coordinate intervals gives
\[
\Delta_V(f)\leq(2\rho)^k\sum_e
 \sup_{y\in\mathcal C_n}|h^{(k)}(a_e\cdot y)|
 \prod_{q\in V}|a_{e,q}|.
\]
For fixed $q$, an elementary symmetric-sum bound yields
\[
\sum_{\substack{|V|=k\\q\in V}}\prod_{j\in V}|a_{e,j}|
 \leq\frac{|a_{e,q}|}{(k-1)!}\left(\sum_j|a_{e,j}|\right)^{k-1}.
\]
Combining this inequality, \eqref{sqt:eq:derivativebounds} and \eqref{sqt:eq:rowsumbounds}, uniformly for $k\leq2n$, proves
\begin{equation}\label{sqt:eq:Dk}
D_1(f),D_2(f)=O(n\rho^3),\qquad
D_k(f)\leq nk(C\rho)^k\quad(k\geq3).
\end{equation}
These bounds include orders above the selected cumulant truncation.

Fix an integer $p\geq1$, and choose
\begin{equation}\label{sqt:eq:cutofforders}
m=4(p+1),\qquad L=11(p+1),\qquad\alpha=n^{-1/2+4\eps}.
\end{equation}
For each fixed $t\leq m$, the left side of \eqref{sqt:eq:isaevhyp} is $O_m(n\rho^3)$. For example, for $t\geq3$ put $z=C\rho e^{3\alpha/2}$ and use
\[
\sum_{k\geq t}k\binom{k-1}{t-1}z^k
 =\frac{tz^t}{(1-z)^{t+1}}.
\]
The first two terms for $t=1,2$ are controlled by \eqref{sqt:eq:Dk}, and the remaining series is again $O_m(n\rho^3)$. As
$\alpha/(n\rho^3)=n^\eps$, all the hypotheses of Proposition~\ref{sqt:prop:complex} hold for sufficiently large $n$, including $\alpha\leq1/100$. The function is bounded, measurable and complex-valued on the required product domain.

The dimension-adjusted error is
\[
2n\alpha^{m+1}=O\left(n^{1-(1/2-4\eps)(4p+5)}\right)=O(n^{-p}).
\]
At $\eps=1/100$ the exponent is $-1.84p-1.30$, leaving ample slack. Thus
\begin{equation}\label{sqt:eq:finitecumulant}
\E e^{f(\widehat X)}
 =(1+O(n^{-p}))\exp\left(\sum_{r=1}^m\frac{\kappa_r(f(\widehat X))}{r!}\right).
\end{equation}
This is where the order-one cubic phase is controlled. Expanding that phase as if it were a uniformly small perturbation would not justify \eqref{sqt:eq:finitecumulant}.

\subsection{Replace only the finite cumulants by polynomials}
Let
\[
R_L(x)=\sum_{d=3}^L h_d\sum_{i,j}(x_i+x_{n+j})^d.
\]
On $\mathcal B_n$,
\begin{equation}\label{sqt:eq:RLbounds}
|R|+|R_L|=O(n^2\rho^3),\qquad |R-R_L|=O(n^2\rho^{L+1}).
\end{equation}
The moment-partition formula for a cumulant, with one factor subtracted at a time, bounds the difference of each order-$r$ cumulant by
\[
C_r\|R-R_L\|_\infty
 (1+\|R\|_\infty+\|R_L\|_\infty)^{r-1}.
\]
For $r\leq m$ the exponent of $n$ in this bound is at most
\begin{equation}\label{sqt:eq:replacementexponent}
1+\frac m2-\frac L2+\eps(L+3m-2)
 \leq-\frac{327}{100}p-\frac{229}{100}<-p.
\end{equation}
Hence replacing $R$ by $R_L$ inside the finite cumulant sum costs $O(n^{-p})$.

For fixed $r,L$, the polynomial $R_L(TX)^r$ has fixed degree; the number of its terms and the magnitude of its coefficients are bounded by powers of $n$. Fixed Gaussian moments are likewise polynomially bounded. By Cauchy--Schwarz and \eqref{sqt:eq:tailprob}, conditioning on $\mathcal C_n$ changes any required moment by at most a polynomial in $n$ times $e^{-c n^{2\eps}}$. Applying the moment-partition formula again gives
\begin{equation}\label{sqt:eq:fullgaussreplace}
\kappa_r(R_L(T\widehat X))=\kappa_r(R_L(TX))+O(n^{-p}),\qquad r\leq m.
\end{equation}
No unbounded exponential of $R_L$ has been introduced in this step.

Set $Y=TX$ and $H_{p,n}=\sum_{r=1}^m\kappa_r(R_L(Y))/r!$. The diagram bound in the next section proves that $H_{p,n}$ is real and $O(1)$. Therefore $e^{H_{p,n}}$ is bounded above and away from zero. Combining \eqref{sqt:eq:transfer}, \eqref{sqt:eq:prefactor}, \eqref{sqt:eq:boundedexpectation}, and \eqref{sqt:eq:finitecumulant}--\eqref{sqt:eq:fullgaussreplace} gives
\begin{equation}\label{sqt:eq:analyticreduction}
a_n=G_n\exp(H_{p,n})(1+O(n^{-p})).
\end{equation}
The exponentially small absolute localization error is now a relative negligible error. Positivity and the real logarithm in subsequent formulas are justified by the real bounded exponent and a factor tending to $1$; no nonvanishing assumption has been added to Isaev's theorem.

\section{Connected contractions and finite power counting}\label{sec:diagrams}
\subsection{The exact covariance and connected Wick identity}
Put $Z_{ij}=Y_i+Y_{n+j}$ and $S_d=\sum_{i,j}Z_{ij}^d$. Since $TT^{\mathsf T}=A^{-1}=I-\mathbf J_{2n}/(4n)$,
\begin{equation}\label{eq:covariance}
\Cov(Y_a,Y_b)=\frac{\delta_{ab}-1/(4n)}{2n},\qquad
\Cov(Z_{ij},Z_{k\ell})=\frac{n\delta_{ik}+n\delta_{j\ell}-1}{2n^2}.
\end{equation}
For a nonempty degree tuple $\boldsymbol d=(d_1,\ldots,d_r)$, $r\geq1$, with every $d_a\geq3$, let
\[
K_{\boldsymbol d}(n)=\kappa(S_{d_1},\ldots,S_{d_r}),\qquad
D=\sum_a d_a.
\]
It vanishes if $D$ is odd. For even $D$, set $E=D/2$ and
\begin{equation}\label{eq:cost}
c(\boldsymbol d)=E-r=\frac12\sum_a(d_a-2).
\end{equation}
This half-sum is the cost throughout the article.

The connected-Wick toolkit and general cost-counting viewpoint are standard; see, for example, Isaev--McKay \cite[Theorems~2.2 and~2.4]{IM}. We give the needed specialized argument here. Wick's formula pairs the $D$ marked half-edges attached to the $r$ labelled polynomial vertices. In a joint cumulant, only pairings whose induced multigraph is connected remain; loops are allowed and a one-vertex graph is connected. For completeness, insert Wick's moment formula into
\begin{equation}\label{eq:cumulantpartition}
\kappa(W_1,\ldots,W_r)=
\sum_{\pi\in\Part([r])}(-1)^{|\pi|-1}(|\pi|-1)!
 \prod_{B\in\pi}\E\prod_{a\in B}W_a.
\end{equation}
For a fixed pairing, each component must lie within one block of $\pi$. Summing the displayed coefficients over partitions of the set of components gives $1$ when there is one component and $0$ otherwise. This follows either from the moment-cumulant identity for a constant variable or the coefficient identity $\log(e^z)=z$. Thus the connected rule follows directly from ordinary Wick and partition identities.

For each paired edge, expand \eqref{eq:covariance} into a row equality term $\delta_{ik}/(2n)$, a column equality term $\delta_{j\ell}/(2n)$, and a constant term $-1/(2n^2)$. Let their counts be $E_R,E_C,E_0$. Let $c_R,c_C$ be the numbers of components, including isolated vertices, in the row-colored and column-colored graphs. Summing freely over all cell indices gives exactly
\begin{equation}\label{eq:colorweight}
2^{-E}(-1)^{E_0}n^{c_R+c_C+E_R+E_C-2E}.
\end{equation}
Loops impose no equality on different indices; this convention handles them automatically. Every cumulant $K_{\boldsymbol d}(n)$ is therefore a finite rational Laurent polynomial in $n$.

A more efficient enumeration uses edge multiplicities. If $m_{ab}$ edges join distinct vertices $a,b$ and $\ell_a$ loops sit at $a$, the number of marked half-edge pairings inducing that multigraph is
\begin{equation}\label{eq:pairmultiplicity}
N(G)=\frac{\prod_a d_a!}
 {\left(\prod_{a<b}m_{ab}!\right)\left(\prod_a2^{\ell_a}\ell_a!\right)}.
\end{equation}
Indeed, assign the marked half-edges at each vertex to its neighbors and loops; match half-edges between each pair of distinct vertices; then pair the $2\ell_a$ loop half-edges. The resulting factorials simplify to \eqref{eq:pairmultiplicity}. Using this multiplicity is exactly equivalent to enumerating marked pairings.

\subsection{A bound before cancellation}
\begin{lemma}\label{lem:power}
Every colored contribution to $K_{\boldsymbol d}(n)$ has Laurent degree between $-2c(\boldsymbol d)$ and $1-c(\boldsymbol d)$. In particular,
\begin{equation}\label{eq:powerbound}
K_{\boldsymbol d}(n)=O\bigl(n^{1-c(\boldsymbol d)}\bigr).
\end{equation}
\end{lemma}
\begin{proof}
Write $q=E_R+E_C$ and let $\rho_R=r-c_R$, $\rho_C=r-c_C$ be the two graphic ranks. Since each nonconstant edge increases rank by at most one,
$\rho_R+\rho_C\leq q$. The exponent in \eqref{eq:colorweight} is
\[
2r-\rho_R-\rho_C+q-2E
 =-2c+(q-\rho_R-\rho_C)\geq-2c.
\]
For the upper bound, let $c_U$ be the number of components in the union of the row- and column-colored graphs. Connectivity of the original pairing graph after adding $E_0$ constant edges implies $c_U\leq E_0+1$. The sum of graphic ranks is at least the rank of the union:
\[
(r-c_R)+(r-c_C)\geq r-c_U.
\]
Thus $c_R+c_C\leq r+E_0+1$, and the exponent is at most
$r+E_0+1+E_R+E_C-2E=r+1-E=1-c$.
\end{proof}
This estimate is termwise. It does not depend on cancellation between colors or on an observed pattern in low-order computations.

\subsection{A formula at every fixed order}
Because $d_a\geq3$ and $D$ is even, every nonzero tuple has integer cost $c\geq1$. Lemma~\ref{lem:power} makes every term in $H_{p,n}$ bounded. Moreover $R_L(-Y)=\overline{R_L(Y)}$ and the Gaussian law is symmetric, so its moments and cumulants are real. These facts establish the claim about $H_{p,n}$ used in \eqref{eq:analyticreduction}.

Define
\begin{equation}\label{eq:tupleweight}
w(\boldsymbol d)=\frac{\prod_{a=1}^r h_{d_a}}{\prod_{d\geq3}m_d!},
\qquad m_d=\#\{a:d_a=d\}.
\end{equation}
The denominator accounts for passing from the ordered cumulant sum, which has factor $1/r!$, to nondecreasing degree tuples. The Taylor factorials are already included in the $h_d$; no further degree factorial is inserted into $w$.

If $c(\boldsymbol d)>p$, \eqref{eq:powerbound} is $O(n^{-p})$. Conversely, if $c\leq p$, then $r\leq2p$ and $d_a\leq2p+2$. Hence every such tuple occurs among the cutoffs $m,L$ in \eqref{eq:cutofforders}. Taking the real logarithm in \eqref{eq:analyticreduction}, we obtain
\begin{equation}\label{eq:finiteformula}
\log\frac{a_n}{G_n}
 =\sum_{\substack{\boldsymbol d\text{ nondecreasing},\ r\geq1,\ d_a\geq3\\
                  D\text{ even},\ c(\boldsymbol d)\leq p}}
 w(\boldsymbol d)K_{\boldsymbol d}(n)+O(n^{-p}).
\end{equation}
Laurent monomials of degree at most $-p$ may be discarded. All retained coefficients are rational: in even total degree the product of powers of $\ii$ in the $h_d$ is real, and the remaining factors are rational.

At cost one only $(4)$ and $(3,3)$ occur. Their constant contributions, computed directly below, are $13/4$ and $-3$, giving $1/4$. With $p=K+1$, equation~\eqref{eq:finiteformula} proves \eqref{eq:main}. Compatibility of the coefficients as $p$ increases follows directly from the degree bound, or from uniqueness of a Poincar\'e expansion.

Here is a finite coefficient algorithm for any prescribed $K$:
\begin{enumerate}
\item For $1\leq c\leq K+1$, list the integer partitions of $2c$ and add $2$ to every part. These are precisely the possible nondecreasing degree tuples.
\item Compute $h_d$ through $d=2K+4$ by \eqref{eq:stirling}.
\item Enumerate connected multigraphs with the prescribed degrees. Weight by \eqref{eq:pairmultiplicity}, then sum \eqref{eq:colorweight} over edge colors.
\item Multiply by \eqref{eq:tupleweight}, add the Laurent polynomials, and retain degrees $0,-1,\ldots,-K$.
\end{enumerate}
This is a finite rational algorithm, not a claim of practical efficiency for large $K$. The public implementation deliberately restricts itself to $K\leq2$; the mathematical algorithm has no such fixed-order restriction.

\section{Two exact correction coefficients}\label{sqt:sec:coefficients}
\subsection{The seven cumulants through the first correction}
The costs one and two give the following seven exact Laurent identities:
\begin{align}
K_{(4)}&=\frac{3(2n-1)^2}{4n^2},&
K_{(6)}&=\frac{15(2n-1)^3}{8n^4},\label{sqt:eq:k46}\\
K_{(3,3)}&=\frac{3(16n^2-12n+1)}{8n^2},&
K_{(3,5)}&=\frac{15(2n-1)(20n^2-12n-1)}{16n^4},\label{sqt:eq:k335}\\
K_{(4,4)}&=\frac{3(13n^3-15n^2+3n+1)}{n^4},\label{sqt:eq:k44}\\
K_{(3,3,4)}&=\frac{9(124n^3-116n^2+10n+7)}{8n^4},\label{sqt:eq:k334}\\
K_{(3,3,3,3)}&=\frac{243(22n^3-18n^2+1)}{8n^4}.\label{sqt:eq:k3333}
\end{align}
They follow from the finite contractions above; the package provides exact independent checks. For example, $\Var(Z_{ij})=(2n-1)/(2n^2)$ and Wick's fourth moment immediately gives $K_{(4)}$. For two centered Gaussian variables $U,V$, Wick gives
\[
\E U^3V^3=9\Var(U)\Var(V)\Cov(U,V)+6\Cov(U,V)^3.
\]
Substituting \eqref{sqt:eq:covariance} and summing the indices gives $K_{(3,3)}$.

In the order $(4),(6),(3,3),(3,5),(4,4),(3,3,4),(3,3,3,3)$, the respective weights are
\[
\frac{13}{12},\quad-\frac{541}{360},\quad-\frac12,\quad
\frac54,\quad\frac{169}{288},\quad-\frac{13}{24},\quad\frac1{24}.
\]
Their weighted sum is
\begin{equation}\label{sqt:eq:lowercostsum}
\frac14-\frac{3}{2n}+\frac{35}{8n^2}+O(n^{-3}).
\end{equation}
In particular the constant is $13/4-\tfrac12\cdot6=1/4$. The cubic sum has limiting variance $6$; it is not asymptotically negligible. This is precisely the order-one phase that the complex cumulant theorem resums.

\subsection{The eleven additional cost three terms}
Through $n^{-2}$, costs $1,2,3$ give $2+5+11=18$ degree multisets. For cost three, the leading Laurent degree is $-2$. Table~\ref{sqt:tab:costthree} supplies that coefficient and the corresponding log-series weight. Multiplying and summing its rows gives $83/32$. Combining with the $35/8$ already present in \eqref{sqt:eq:lowercostsum} gives
\[
C_2=\frac{35}{8}+\frac{83}{32}=\frac{223}{32}.
\]
Exponentiating adds $C_1^2/2=9/8$ to the second relative coefficient, yielding $259/32$ in \eqref{sqt:eq:relativetwo}.

\begin{table}[htbp]
\centering
\caption{All new terms at cost three. The middle column is the coefficient of $n^{-2}$ in the exact cumulant.}\label{sqt:tab:costthree}
\begin{tabular}{lrr}
\toprule
Degree tuple $\boldsymbol d$ & $[n^{-2}]K_{\boldsymbol d}$ & $w(\boldsymbol d)$\\
\midrule
$(8)$ & $105$ & $47293/20160$\\
$(3,7)$ & $315$ & $-223/120$\\
$(4,6)$ & $315$ & $-7033/4320$\\
$(5,5)$ & $270$ & $-25/32$\\
$(3,3,6)$ & $2655/2$ & $541/720$\\
$(3,4,5)$ & $2475/2$ & $65/48$\\
$(4,4,4)$ & $1350$ & $2197/10368$\\
$(3,3,3,5)$ & $6885$ & $-5/24$\\
$(3,3,4,4)$ & $6813$ & $-169/576$\\
$(3,3,3,3,4)$ & $89181/2$ & $13/288$\\
$(3,3,3,3,3,3)$ & $681615/2$ & $-1/720$\\
\bottomrule
\end{tabular}
\end{table}
Ignoring the lower-cost $n^{-2}$ terms would give the incorrect answer $83/32$. Omitting the six-cubic cumulant would also change the result substantially. The graph calculation includes all $3550$ labelled connected graphs with six cubic vertices, with loops and multiple edges allowed.

\subsection{An independent formal factorization}
The second verification computes all eighteen \emph{full} Laurent polynomials without generating connected graphs, coloring edges, fitting numerical values, or interpolating. Introduce a formal Gaussian Wick functional under which $A_1,\ldots,A_n,B_1,\ldots,B_n,C$ are independent and
\[
\E A_i^2=\E B_j^2=1,\qquad\E C^2=-\frac1n.
\]
This is an algebraic moment functional, not a probability distribution. It assigns zero to odd Gaussian moments and the usual pairing value to even ones, allowing the displayed negative covariance. Then
\[
Z_{ij}=\frac{A_i+B_j+C}{\sqrt{2n}}
\]
has exactly the covariance \eqref{sqt:eq:covariance}. Since Gaussian polynomial moments are determined by their pair covariances, the formal functional reproduces the required true-Gaussian polynomial moments.

Let $P_a(A)=\sum_iA_i^a$ with $P_0(A)=n$, and similarly for $B$. The multinomial identity gives
\begin{equation}\label{sqt:eq:powerfactor}
S_d=(2n)^{-d/2}\sum_{a+b+c=d}\frac{d!}{a!b!c!}
 P_a(A)P_b(B)C^c.
\end{equation}
For positive exponents $u_1,\ldots,u_q$, decomposing the chosen row indices by their equality partition gives
\begin{equation}\label{sqt:eq:powersummoments}
\E\prod_{j=1}^qP_{u_j}(A)
 =\sum_{\pi\in\Part([q])}(n)_{|\pi|}
   \prod_{B\in\pi}\E G^{\sum_{j\in B}u_j},
\end{equation}
where $G$ is standard normal and $(n)_k=n(n-1)\cdots(n-k+1)$. Factors $P_0$ are first replaced by $n$. The $B$ moments have the same formula, independently, and
\[
\E C^c=\begin{cases}0,&c\text{ odd},\\(c-1)!!(-1/n)^{c/2},&c\text{ even}.
\end{cases}
\]
Expanding products of \eqref{sqt:eq:powerfactor}, using \eqref{sqt:eq:powersummoments}, and finally applying \eqref{sqt:eq:cumulantpartition} gives a separate exact derivation. It agrees identically with all eighteen full Laurent polynomials, not just the three coefficients retained here.

Additional direct checks at $n=1,2$ use actual Gaussian variables. At $n=2$ the four cells can be represented as $\pm a\pm b+w$, with independent $a,b,w\sim N(0,1/8)$ and the row and column signs chosen separately. This reproduces \eqref{sqt:eq:covariance} and avoids the formal negative variance. These algebraic cross-checks support the coefficient computation; the analytic remainder is supplied by Sections~\ref{sqt:sec:gauge}--\ref{sqt:sec:diagrams}.

\section{Controlled inversion and logarithmic reversion}\label{sec:inverse}
\subsection{Strict monotonicity of the count}
For any $n\times n$ table counted by $a_n$, add $1$ to its $n$ diagonal entries, append a zero last row and last column, and place $n+1$ in the new bottom-right entry. This gives an injective map into the tables counted by $a_{n+1}$. The all-ones $(n+1)\times(n+1)$ table is not in its image, because its new off-diagonal entries are nonzero. Thus $a_{n+1}>a_n$ for $n\geq1$. In particular \eqref{eq:inversecount} is a well-defined integer threshold.

\subsection{An asymptotic two-ceiling enclosure}
For fixed $K$, differentiation of \eqref{eq:FK} gives
\begin{equation}\label{eq:FKprime}
F_K'(x)=2x\log4-\log x-\log(4\pi)-1+\frac1x+O_K(x^{-2}).
\end{equation}
It is positive and comparable to $x$ on a sufficiently large real half-line. Since $F_K(x)\to\infty$, there is a unique root $t_K$ on this increasing branch for sufficiently large $Y$.

\begin{theorem}\label{thm:inverse}
For each fixed $K\geq0$, there exist positive constants $B_K,D_K,N_K$ such that, for sufficiently large $Y$, the roots $x_-,x_+$ on the increasing branches defined by
\begin{align}
F_K(x_-)+B_Kx_-^{-K-1}&=\log Y,\label{eq:minusroot}\\
F_K(x_+)-B_Kx_+^{-K-1}&=\log Y\label{eq:plusroot}
\end{align}
satisfy
\begin{equation}\label{eq:tworoots}
\ceil{x_-}\leq\nu(Y)\leq\ceil{x_+},\qquad
x_+-t_K=O_K(t_K^{-K-2}),\quad t_K-x_-=O_K(t_K^{-K-2}).
\end{equation}
In particular \eqref{eq:inverseintro} holds. At $K=2$, the real-index uncertainty is $O(t_2^{-4})$.
\end{theorem}
\begin{proof}
Theorem~\ref{thm:main} gives
$|\log a_n-F_K(n)|\leq B_Kn^{-K-1}$ for $n\geq N_K$. Both perturbed real functions in \eqref{eq:minusroot}--\eqref{eq:plusroot} have derivatives comparable to $x$ for large $x$. If an integer $n<x_-$ lies in that range, its upper log bound is less than $\log Y$, so $a_n<Y$. If $n\geq x_+$, its lower log bound is at least $\log Y$, so $a_n\geq Y$. Increasing $Y$ disposes of the finitely many smaller integers and proves the ceiling inequalities. At $t_K$, the two perturbations are $\pm B_Kt_K^{-K-1}$; the mean value theorem and \eqref{eq:FKprime} place the roots within $O_K(t_K^{-K-2})$. Enlarging the constant gives \eqref{eq:inverseintro}.
\end{proof}
The constants in this theorem are existential asymptotic constants. Computing a root of $F_K$ to many decimal places does not compute $B_K$ or prove a certified finite-$Y$ enclosure. Near an integer, the two ceilings can differ, even when their real arguments are very close.

\subsection{Four explicit terms of the inverse}
Write
\[
a=\log4,\qquad b=\log(4\pi),\qquad
r=\sqrt{\frac{\log Y}{a}},\qquad \ell=\log r.
\]
Define polynomials in $\ell$, with constants depending on $a,b$, by
\begin{align}
d&=\frac{\ell+b}{2a},\label{eq:reversiond}\\
e&=\frac{ad^2+d-\ell-b/2-1/4}{2a},\label{eq:reversione}\\
f&=\frac{e-d+d^2/2-C_1}{2a},\label{eq:reversionf}\\
g&=\frac{f-ae^2+de-d^3/6-e+d^2/2+C_1d-C_2}{2a},\label{eq:reversiong}
\end{align}
where $C_1=-3/2$ and $C_2=223/32$. Then the implicit root $t_2$ satisfies
\begin{equation}\label{eq:reversion}
t_2=r+d+\frac e r+\frac f{r^2}+\frac g{r^3}
       +O\left(\frac{(\log r)^5}{r^4}\right).
\end{equation}
Here the displayed logarithmic remainder is deliberately conservative. It is a truncation error for the explicit log-polynomial expression, distinct from the $O(t_2^{-4})$ uncertainty between the implicit asymptotic threshold and the sequence.

To check \eqref{eq:reversion}, put $s=d+e/r+f/r^2+g/r^3$ and use
\begin{align*}
F_2(r+s)-ar^2={}&r(2as-\ell-b)
 +as^2-s(\ell+b)-s+\ell+b/2+1/4\\
&+\frac{-s^2/2+s+C_1}{r}
 +\frac{s^3/6-s^2/2-C_1s+C_2}{r^2}
 +O\left(\frac{(\log r)^5}{r^3}\right).
\end{align*}
The coefficients of $r,1,r^{-1},r^{-2}$ vanish successively by \eqref{eq:reversiond}--\eqref{eq:reversiong}, using $2ad=\ell+b$. Since $d=O(\ell)$ and $e,f,g=O(\ell^4)$ collectively, ordinary Taylor bounds for $\log(1+s/r)$ and $(1+s/r)^{-1}$ justify the residual estimate. The derivative is comparable to $r$, so the mean value theorem gives \eqref{eq:reversion}. The companion exact symbolic check verifies these four cancellations treating $a,b,\ell,C_1,C_2$ as indeterminates.

\subsection{Finite-step inversion at any fixed order}
For any fixed $K$, start Newton iteration at $x_0=r+d$:
\begin{equation}\label{eq:newton}
x_{j+1}=x_j-\frac{F_K(x_j)-\log Y}{F_K'(x_j)}.
\end{equation}
The first two reversion terms imply $x_0-t_K=O((\log r)^2/r)$. In a fixed relative neighborhood of $r$, $F_K'\asymp r$ and $F_K''=O_K(1)$. Taylor's theorem for one Newton step therefore gives
\[
|x_{j+1}-t_K|\leq\frac{C_K}{r}|x_j-t_K|^2.
\]
By induction, for each fixed $j$,
\begin{equation}\label{eq:newtonerror}
x_j-t_K=O_K\left(
 r^{-(2^{j+1}-1)}(\log r)^{2^{j+1}}\right).
\end{equation}
Choosing $2^{j+1}-1>K+2$ makes this numerical-iteration approximation asymptotically smaller than the sequence uncertainty in Theorem~\ref{thm:inverse}. For $K=2$, two Newton steps from $r+d$ suffice. This statement concerns exact-real iteration; floating-point evaluation must additionally control its own rounding error.

Equations~\eqref{eq:reversion} and \eqref{eq:newtonerror} provide fixed-order logarithmic reversion. If one calls such expansions the algebraic part of a formal transseries, the qualification matters: no nonperturbative exponential sectors, resurgence statement, or Stokes behavior have been established here.

\section{A compact positive density extension}\label{sqt:sec:density}
The square proof extends with uniform constants on each fixed compact interval of positive densities. This section specifies the quantifiers and all changes to the analytic argument; it does not assert uniformity in sparse or unbounded density regimes.

Let $T(n,s)$ count labelled $n\times n$ nonnegative integer tables with all line sums $s$, and write
\begin{equation}\label{sqt:eq:densitydefs}
\lambda=\frac{s}{n},\qquad u=\lambda(1+\lambda),\qquad A_\lambda=\frac u2,
\qquad H(\lambda)=\frac{(1+\lambda)^{1+\lambda}}{\lambda^\lambda}.
\end{equation}
The matrix $A$ of \eqref{sqt:eq:augmented} is unrelated to the scalar $A_\lambda$; the whitening map $T$ stays as in \eqref{sqt:eq:whitening}. Define
\begin{equation}\label{sqt:eq:densityG}
\mathcal G(n,\lambda)=\frac{H(\lambda)^{n^2}}
 {(2\pi u)^{n-1/2}n^{n-1}},\qquad
B_0(u)=\frac13-\frac1{6u}.
\end{equation}

\begin{theorem}\label{sqt:thm:density}
There are polynomials $P_j\in\mathbb Q[z]$, of degree at most $j+1$, such that for every fixed compact interval $I=[\lambda_-,\lambda_+]\subset(0,\infty)$ and every fixed integer $K\geq0$,
\begin{equation}\label{sqt:eq:densitytheorem}
\log\frac{T(n,s)}{\mathcal G(n,s/n)}
 =B_0(u)+\sum_{j=1}^K\frac{P_j(1/u)}{n^j}
   +O_{I,K}(n^{-K-1}),
\end{equation}
uniformly over positive integers $n,s$ with $s/n\in I$. In particular,
\begin{equation}\label{sqt:eq:densitypolys}
P_1(z)=-\frac32,\qquad
P_2(z)=\frac{1171}{180}+\frac{11}{12}z+\frac1{60}z^2+\frac1{180}z^3.
\end{equation}
At $\lambda=1$, this recovers Theorem~\ref{sqt:thm:main}.
\end{theorem}
The statement permits varying admissible densities $s_n/n$. For a fixed rational density, it applies along the admissible subsequence. At a fixed irrational density, exact integral margins are unavailable; rounding $\lambda n$ changes the density and must be reflected in every occurrence of $\lambda$ in \eqref{sqt:eq:densitytheorem}.

\subsection{Uniform localization and the scaled gauge}
The centered characteristic function is now
\[
\psi_\lambda(z)=\frac{e^{-\ii\lambda z}}{1-\lambda(e^{\ii z}-1)},\qquad
\mathcal F_\lambda(\theta,\phi)=\prod_{i,j}\psi_\lambda(\theta_i+\phi_j).
\]
Although an individual factor need not be $2\pi$-periodic, the full product is periodic in each angle because $n\lambda=s\in\mathbb Z$. Gauge invariance is exact. Each table has geometric product probability $H(\lambda)^{-n^2}$, giving
\[
T(n,s)=H(\lambda)^{n^2}(2\pi)^{-2n}\int\mathcal F_\lambda.
\]

Canfield--McKay's condition~(1.1) specializes to
\[
\frac23\left(4+\frac1u\right)\leq a\log n.
\]
With fixed $a=b=1/8$, it holds uniformly on $I$ for all sufficiently large $n$; their polynomial density bound is automatic. Their stated asymptotic constants allow dependence on $a,b$ and the fixed localization exponent, with margins among the asymptotic variables. One can also check uniformity directly in their localization proof. The subdivision count $\lceil6000(1+\lambda)\rceil$ is bounded on $I$, its angular width is bounded above and below, and $A_\lambda$, $A_\lambda^{-1}$ and $\lambda/(1+\lambda)$ have the needed compact bounds. The far-region decay is $\exp\{-c_I n^{1+\eta}\}$ against a reciprocal Gaussian normalization $\exp\{O_I(n\log n)\}$. The exceptional-coordinate decay is $\exp\{-c_I n(m_2+n_2)\}$, and the remaining omitted small intervals contribute $\exp\{-c_I n^{2\eta}\}$. Their Lemma~4 is uniform in $\lambda>0$. More precisely, after normalization the three bounds have the forms
\[
 e^{C_I n\log n-c_I n^{1+\eta}},\qquad
 \sum_{1\leq r\leq2n^\eta}n^{C_Ir+C_I}e^{-c_Inr},\qquad
 n^{C_I}e^{-c_In^{2\eta}},
\]
where $r=m_2+n_2$ counts exceptional coordinates. The middle sum decays geometrically in $r$ for sufficiently large $n$. Each expression is $O_I(e^{-c'_In^\eta})$ for a sufficiently small fixed $\eta$. Thus
\begin{equation}\label{sqt:eq:densityminor}
\int_{\mathcal R_{\lambda,\eta}^c}|\mathcal F_\lambda|
 =O_I(e^{-c_I n^\eta})I_{0,\lambda},\qquad
I_{0,\lambda}=\pi^{n-1/2}A_\lambda^{-n+1/2}n^{1-n}
 e^{-(1+2A_\lambda)/(12A_\lambda)}.
\end{equation}
The region has centered deviations bounded by $(1+\lambda)^{-1}n^{-1/2+\eta}$ and mean sum bounded by $(1+\lambda)^{-1}n^{-1+4\eta}$. These bounds make its small $v^\perp$ representative $o_I(\rho)$ uniformly, as before.

The correct augmented weight is $\exp\{-A_\lambda(v\cdot x)^2/2\}$. Its quadratic part is $-A_\lambda q_+(x)$ and its full Gaussian integral is
\[
Z_{n,\lambda}=\frac1{\sqrt2}\left(\frac\pi{A_\lambda n}\right)^n.
\]
The real gauge fiber has mass $\sqrt{2\pi/A_\lambda}/(2n)$, so the conversion factor is $Q_{n,\lambda}=2n\sqrt{2\pi A_\lambda}$. In particular,
\[
\frac{Q_{n,\lambda}Z_{n,\lambda}}{(2\pi)^{2n}}
 =(4\pi A_\lambda)^{-n+1/2}n^{1-n},\qquad
\frac{I_{0,\lambda}}{Q_{n,\lambda}Z_{n,\lambda}}
 =\frac{e^{-(1+2A_\lambda)/(12A_\lambda)}}{2\pi}.
\]
The last ratio is uniformly bounded above and below on $I$. The proof of Lemma~\ref{sqt:lem:transfer} now applies verbatim with this weight: the omitted gauge tail is $e^{-\Omega_I(n^{1+2\eps})}$ and its normalization costs at most $e^{O_I(n\log n)}$. This proves a uniform torus-to-cube transfer. Using an unscaled gauge while retaining the scaled Gaussian covariance would give an incorrect normalization.

\subsection{Uniform analytic and cumulant estimates}
Define
\[
h_\lambda(z)=-\ii\lambda z-\log[1-\lambda(e^{\ii z}-1)]+A_\lambda z^2
 =\sum_{d\geq3}h_d(\lambda)z^d.
\]
The nearest singularity is at $-\ii\log(1+1/\lambda)$. For example, $r_0=\tfrac14\log(1+1/\lambda_+)>0$ gives a uniform analytic disk of radius $2r_0$, because
$\lambda_+(e^{2r_0}-1)<1$. The logarithmic series on this disk is uniformly bounded. Hence \eqref{sqt:eq:derivativebounds} holds with constants depending only on $I$, and
\begin{equation}\label{sqt:eq:densitycumulants}
h_d(\lambda)=\frac{\ii^d}{d!}\kappa_d(\lambda),\qquad
\kappa_d(\lambda)=\sum_{j=1}^d\genfrac{\{}{\}}{0pt}{}{d}{j}(j-1)!\lambda^j.
\end{equation}
The independent Gaussian coordinates before applying $T$ now have variance $1/(2A_\lambda n)$. Conditioning on the same cube preserves independence, with complement probability $O_I(e^{-c_I n^{2\eps}})$. The coefficient-vector bounds \eqref{sqt:eq:rowsumbounds} are unchanged. Therefore \eqref{sqt:eq:Dk} holds uniformly, and the same $\alpha,m,L$ in \eqref{sqt:eq:cutofforders} satisfy Isaev's hypotheses uniformly for sufficiently large $n$.

Taylor replacement inside the finite cumulants has the same exponent \eqref{sqt:eq:replacementexponent}; its constants depend on $I,p$. Fixed-degree Gaussian moments remain uniformly polynomially bounded since $A_\lambda$ stays away from zero. Removing the conditioning again costs less than every inverse power of $n$. The resulting finite exponent is real and uniformly bounded by the connected-diagram estimate below, so absolute localization errors are relative negligible errors. This establishes the uniform version of \eqref{sqt:eq:analyticreduction}.

\subsection{The density polynomials and their coefficients}
The pair-sum covariance becomes
\[
\Cov(Z_{ij},Z_{k\ell})=\frac{n\delta_{ik}+n\delta_{j\ell}-1}{2A_\lambda n^2}.
\]
For a tuple of even total degree $2E$, its Gaussian cumulant is exactly $A_\lambda^{-E}K_{\boldsymbol d}(n)$, with $K_{\boldsymbol d}$ the density-one polynomial. Hence the uniform finite formula is
\begin{equation}\label{sqt:eq:densityfinite}
\log\frac{T(n,s)}{\mathcal G(n,\lambda)}
 =\sum_{\substack{\boldsymbol d\text{ nondecreasing},\ r\geq1,\ d_a\geq3\\
                    D\text{ even},\ c(\boldsymbol d)\leq p}}
 \frac{\prod_a h_{d_a}(\lambda)}{\prod_d m_d!}
 A_\lambda^{-E}K_{\boldsymbol d}(n)+O_{I,p}(n^{-p}).
\end{equation}
The power bound is unchanged. This already proves existence, uniformity and a finite rational-function coefficient algorithm.

To obtain the polynomial structure, use
\[
\kappa_2=u,\qquad \kappa_{d+1}=u\frac{\dd}{\dd\lambda}\kappa_d.
\]
Inductively, $\kappa_d$ is divisible by $u$, has degree $d$, and satisfies
$\kappa_d(-1-\lambda)=(-1)^d\kappa_d(\lambda)$. For an even-total-degree tuple,
$\prod_a\kappa_{d_a}/u^r$ is a polynomial invariant under $\lambda\mapsto-1-\lambda$ of degree at most $2(E-r)$. Any such polynomial is a polynomial in $u$ of degree at most $E-r$: set $w=1+2\lambda$ and note that invariance means an even polynomial in $w$, with $w^2=1+4u$. Multiplication by $A_\lambda^{-E}=(2/u)^E$ therefore makes each weighted tuple a polynomial in $1/u$ of degree at most $E-r=c$. To contribute to $n^{-j}$ it must have $c\leq j+1$, proving the asserted degree bound.

The cost-one constant is
\[
\frac{1+6u}{2u}-\frac{2(1+4u)}{3u}=B_0(u).
\]
For the first correction, put $w=1+2\lambda$. The seven tuple contributions, in the order $(3,3),(4),(3,3,3,3),(3,3,4),(3,5),(4,4),(6)$, are
\begin{align*}
&\frac{w^2}{2u},\quad-\frac{1+6u}{2u},\quad
\frac{11w^4}{8u^2},\quad-\frac{31w^2(1+6u)}{12u^2},\\
&\frac{5w^2(1+12u)}{6u^2},\quad
\frac{13(1+6u)^2}{24u^2},\quad-\frac{1+30u+120u^2}{6u^2}.
\end{align*}
With $w^2=1+4u$, the first two sum to $-1$ and the remaining five to $-1/2$. All eighteen tuples at costs at most three give
\[
P_2(1/u)=\frac{1171u^3+165u^2+3u+1}{180u^3}.
\]
The same independent formal-factorization computation described in Section~\ref{sqt:sec:coefficients}, with a separately expanded geometric logarithm and the scalar factor $A_\lambda^{-E}$, verifies these rational-function identities exactly. Taking $p=K+1$ in \eqref{sqt:eq:densityfinite} completes the proof of Theorem~\ref{sqt:thm:density}.

\subsection*{[write] The uniformity of the imported localization}
\begin{writenote}[Added 6 October 2026, batch~108]
The uniform minor-arc bound~\eqref{sqt:eq:densityminor} is not a theorem
stated by Canfield and McKay: their Theorem~3 is an asymptotic statement for
$m,n\to\infty$ along admissible parameters, and the source supports
uniformity on $I$ by its own reading of their proof (``One can also check
uniformity directly in their localization proof'', the paragraph
before~\eqref{sqt:eq:densityminor}). Lemma~\ref{sqt:lem:sequential} derives
from the \emph{statement} of their Theorem~3 the weaker uniform bound that
the proof of Theorem~\ref{sqt:thm:density} actually uses;
Remark~\ref{sqt:rem:uniformity} reports what the write found in their proof.
\end{writenote}

\begin{lemma}[{[write]} Uniformity from the sequential statement]\label{sqt:lem:sequential}
Fix a compact $I=[\lambda_-,\lambda_+]\subset(0,\infty)$ and $\eta>0$. For
integers $n,s\ge1$ with $\lambda=s/n\in I$ let $\mathcal R_{\lambda,\eta}$ be
Canfield--McKay's region~(3.1) of \cite{CM} at $m=n$, $s=t$, with their
$\varepsilon$ equal to $\eta$, and put
\[
 E_\eta(n,s)=\frac1{I_{0,\lambda}}\int_{\mathcal R_{\lambda,\eta}^c}|\mathcal F_\lambda| .
\]
Then for every $p>0$,
$\sup\{n^pE_\eta(n,s):s\in\mathbb Z_{\ge1},\ s/n\in I\}\to0$ as $n\to\infty$.
The only input is \cite[Theorem~3]{CM}, read as a statement about every
sequence of admissible tuples with $m,n\to\infty$ along which~(1.1) holds and
$\lambda=n^{O(1)}$.
\end{lemma}
\begin{proof}
Suppose not. Then there are $p>0$, $c>0$, integers $n_1<n_2<\cdots$ and
$s_k$ with $s_k/n_k\in I$ and $n_k^pE_\eta(n_k,s_k)\ge c$ for all $k$.
Consider the tuples $(m,s;n,t)=(n_k,s_k;n_k,s_k)$; then $ms=nt$, and
$s=s(m,n)$ is a well-defined function on the pairs $(n_k,n_k)$ because the
$n_k$ are distinct. At $m=n$ the left side of~(1.1) is
$\tfrac23(4+1/u)$ with $u=\lambda(1+\lambda)\ge u_-=\lambda_-(1+\lambda_-)>0$,
so~(1.1) holds with $a=b=1/8$ as soon as
$\log n_k\ge\tfrac{16}3(4+1/u_-)$, and $\lambda_k\le\lambda_+$ gives
$\lambda=n^{O(1)}$. At $m=n$, $s=t$ the integrand $F$ of~(2.3) of
\cite{CM} is $\mathcal F_\lambda$ and their $I_0$ of~(4.1) is $I_{0,\lambda}$
(Section~\ref{sqt:sec:provenance}). Hence Theorem~3 of \cite{CM} applies along
this sequence: there is $\varepsilon'\in(0,\eta]$ (any sufficiently small one
will do) and a constant $C$ with
$\int_{\mathcal R_{\lambda_k,\varepsilon'}^c}|\mathcal F_{\lambda_k}|\le
Ce^{-n_k^{\varepsilon'}}I_{0,\lambda_k}$ for all large $k$. The three bounds
defining~(3.1) at $m=n$, namely $(1+\lambda)^{-1}n^{-1+4\varepsilon}$ and
twice $(1+\lambda)^{-1}n^{-1/2+\varepsilon}$, do not decrease as
$\varepsilon$ increases ($n\ge1$), and the half-circle condition does not
involve $\varepsilon$, so
$\mathcal R_{\lambda,\varepsilon'}\subseteq\mathcal R_{\lambda,\eta}$ and
$E_\eta(n_k,s_k)\le Ce^{-n_k^{\varepsilon'}}$. Then
$n_k^pE_\eta(n_k,s_k)\to0$, a contradiction.
\end{proof}

The proof of Theorem~\ref{sqt:thm:density} uses~\eqref{sqt:eq:densityminor}
only through the transfer Lemma~\ref{sqt:lem:transfer} with the scaled
weight, where the minor-arc integral enters twice (on
$\mathcal B_n\setminus\widetilde D$, after division by the gauge factor, and
outside the region) and is compared with $Q_{n,\lambda}Z_{n,\lambda}$, whose
ratio to $I_{0,\lambda}$ is bounded above and below on~$I$. Every later error
is $O_{I,p}(n^{-p})$, and the main term is of exact order $Z_{n,\lambda}$
because the exponent $H$ is real and uniformly bounded. So
Lemma~\ref{sqt:lem:sequential}, with $\eta=\varepsilon/2$, may replace the
uniform rate $O_I(e^{-c_In^\eta})$ in~\eqref{sqt:eq:densityminor}: the
torus-to-cube transfer then holds uniformly on $I$ with an error
$O_{I,p}(n^{-p})Z_{n,\lambda}$ for every $p$, and
Theorem~\ref{sqt:thm:density} and Corollary~\ref{sqt:cor:delta} follow as
printed. The same argument, with $\lambda=1$ fixed, is not needed for
Theorem~\ref{sqt:thm:main}, where Theorem~3 of \cite{CM} applies directly.

\begin{remark}[{[write]} What Canfield and McKay state, and what the write found in their proof]\label{sqt:rem:uniformity}
Theorem~3 of \cite{CM} assumes that ``$m,n\to\infty$ in such a way that~(1.1)
holds and $\lambda=n^{O(1)}$'' and concludes, ``for sufficiently small
$\varepsilon>0$'', that $\int_{\mathcal R^c}|F|=O(e^{-n^\varepsilon})I_0$;
their convention (p.~4) lets the implied constants depend on $a$, $b$
and~$\varepsilon$. No version uniform in a parameter family is stated, which
is why Lemma~\ref{sqt:lem:sequential} is needed. The write read their proof
of Theorem~3 (pp.~15--20 of the preprint). It bounds three pieces: the
region where at least $\tfrac13\min(mn^\varepsilon,m^\varepsilon n)$ pairs
have $\cos(\theta_j+\phi_k)\le\cos\delta$, by
$(2\pi)^{m+n}\exp\{-c_1(\log n)^{-1}\min(m^\varepsilon n,mn^\varepsilon)\}$
against $I_0=\exp\{O(m\log n+n\log m)\}$; the regions with $m_2+n_2\ge1$
exceptional coordinates, by a sum of
$n^{O(m_2+n_2)}\exp\{-c_3\lambda(1+\lambda)^{-1}(mn_2+nm_2)\}I_0$, giving
$O(e^{-n^{1-\varepsilon}})I_0$; and the rest of the complement of~(3.1), by
$e^{-n^\varepsilon}I_0$ through a factor
$\exp\{-\tfrac15\lambda(1+\lambda)^{-1}\min(m^{2\varepsilon},n^{2\varepsilon})\}$.
These are the three bounds the source displays before~\eqref{sqt:eq:densityminor}.
The parameter $\lambda$ enters only through $\lambda(1+\lambda)^{-1}$
(bounded below by $(\log n)^{-1}$ by~(1.1)), $A=O(\log n)$,
$N=\lceil6000(1+\lambda)\rceil$, $\delta=2\pi/N$ and their Lemmas~3 and~4,
of which Lemma~4 is stated ``uniformly for $\lambda>0$'' and both are stated
with proofs omitted (``We begin with two technical lemmas whose proofs are
omitted''). On a compact $I$ each of these quantities is bounded above and
below, so on the steps it read the write agrees with the source's reading;
but the implied constants (the $O(\cdot)$ exponents in
$n^{O(m_2+n_2)}$, the factor $O(\delta)$, the constants of Lemma~4) are not
tracked in their text, so the uniform \emph{rate}
$O_I(e^{-c_In^\eta})$ of~\eqref{sqt:eq:densityminor} remains the source's
reading of their proof, not a quoted theorem
(Question~\ref{sqt:q:uniform}). Nothing in Theorem~\ref{sqt:thm:density}
depends on that rate (Lemma~\ref{sqt:lem:sequential}). The same holds for
the source's ``The proof of Lemma~\ref{sqt:lem:transfer} now applies
verbatim with this weight'': it is the source's adaptation, whose two
normalization identities the write rechecked.
\end{remark}

\section{The correction-factor conjecture in an eventual dense square regime}\label{sqt:sec:conjecture}
Canfield--McKay's Conjecture~1 \cite[Eq.~(1.4)]{CM} defines a correction $\Delta(m,s;n,t)$ by
\begin{align}\label{sqt:eq:CMdelta}
M(m,s;n,t)={}&G_{\rm Good}(m,s;n,t)
 \left(1+\frac1m\right)^{(m-1)/2}
 \left(1+\frac1n\right)^{(n-1)/2}\notag\\
&\times\exp\left(-\frac12+\frac{\Delta(m,s;n,t)}{m+n}\right),
\end{align}
where $ms=nt$, $M$ is their table count, and
\[
G_{\rm Good}(m,s;n,t)=
\frac{\binom{n+s-1}{s}^{m}\binom{m+t-1}{t}^{n}}
 {\binom{mn+ms-1}{ms}}.
\]
They conjecture $0<\Delta<2$ for all positive integer admissible parameters. This $G_{\rm Good}$ differs from the Gaussian prefactor $G_n$ of \eqref{sqt:eq:leading}.

\begin{corollary}\label{sqt:cor:delta}
On each fixed compact interval $I\subset(0,\infty)$, uniformly over square tables with integer line sum $s$ and $\lambda=s/n\in I$,
\begin{equation}\label{sqt:eq:deltaasymp}
\Delta(n,s;n,s)=\frac1n\left(\frac{67}{6}+\frac5{3u}\right)+O_I(n^{-2}),
\qquad u=\lambda(1+\lambda).
\end{equation}
Consequently $0<\Delta(n,s;n,s)<2$ for all sufficiently large $n$, uniformly over those margins. In particular,
\begin{equation}\label{sqt:eq:deltadensityone}
\Delta(n,n;n,n)=\frac{12}{n}+O(n^{-2}).
\end{equation}
\end{corollary}
\begin{proof}
Set $\mathcal L(n,\lambda)=\mathcal G(n,\lambda)e^{B_0(u)}$. The standard Stirling expansion of the three factorial ratios in $G_{\rm Good}$ gives, uniformly on $I$,
\begin{equation}\label{sqt:eq:goodstirling}
\log\frac{G_{\rm Good}(n,s;n,s)e^{1/2}}{\mathcal L(n,\lambda)}
 =\frac{g_2(u)}{n^2}+O_I(n^{-4}),
\end{equation}
where
\begin{equation}\label{sqt:eq:g2}
g_2(u)=\frac4{45}+\frac1{12u}+\frac1{60u^2}+\frac1{180u^3}.
\end{equation}
Here is an explicit way to see the coefficient. If
$B(N,\lambda)=\binom{(1+\lambda)N-1}{\lambda N}$, write
\[
\log B(N,\lambda)=N\log H(\lambda)-\tfrac12\log(2\pi N u)
 +\frac{b_1(u)}N+\frac{b_3(u)}{N^3}+O_I(N^{-5}),
\]
with
\[
b_1(u)=-\frac{1+1/u}{12},\qquad
b_3(u)=\frac{1+3/u^2+1/u^3}{360}.
\]
The constant in this display needs the factor $1/(1+\lambda)$ in
$\binom{(1+\lambda)N-1}{\lambda N}=\binom{(1+\lambda)N}{\lambda N}/(1+\lambda)$; it is already included in $-\tfrac12\log(2\pi N u)$. Since
$\log G_{\rm Good}=2n\log B(n,\lambda)-\log B(n^2,\lambda)$,
its $n^{-2}$ coefficient is $2b_3-b_1=g_2$. Its constant is $2b_1=B_0-1/2$. This proves \eqref{sqt:eq:goodstirling}.

Taking logarithms in \eqref{sqt:eq:CMdelta} in the square case,
\[
\frac{\Delta}{2n}=\log\frac{T(n,s)}{G_{\rm Good}}-(n-1)\log(1+1/n)+\frac12.
\]
Use Theorem~\ref{sqt:thm:density}, \eqref{sqt:eq:goodstirling}, and
\[
(n-1)\log(1+1/n)=1-\frac{3}{2n}+\frac5{6n^2}+O(n^{-3}).
\]
The constants and $n^{-1}$ terms cancel. The remaining coefficient is
\[
P_2(1/u)-g_2(u)-\frac56=\frac{67}{12}+\frac5{6u}.
\]
Multiplication by $2n$ proves \eqref{sqt:eq:deltaasymp}. The leading coefficient is positive and bounded above and below on $I$, so its uniform remainder proves the eventual strict inequalities. Substituting $u=2$ gives \eqref{sqt:eq:deltadensityone}.
\end{proof}

This verifies the conjectured inequalities in one regime only: square tables, density in a fixed compact subset of $(0,\infty)$, and $n$ beyond a threshold that is not made explicit (Remark~\ref{sqt:rem:cmregime}). It does not settle all finite sizes, nonsquare tables, or densities tending to zero or infinity. Prior progress includes Greenhill--McKay \cite[Corollary~4.2]{GM}, which proves
$\Delta(m,s;n,t)\sim5(s+t)/(6st)$ in the sparse regime
$st=o((mn)^{1/5})$. Their preceding Corollary~4.1 also contains $-3/(4n)-3/(4m)$, hence $-3/(2n)$ for squares, under sparse assumptions. Those assumptions exclude fixed positive square density. Thus the numerical first coefficient has a sparse-regime antecedent; the present assertion concerns its analytically controlled dense regime and the next correction.

The explicit prefactor in the original conjecture also suggests a term $-3/(2n)$ on expansion. It did not prove that coefficient, because the unknown term $\Delta/(2n)$ could have contributed at the same order. The present corollary shows $\Delta\to0$ in the stated regime.

\begin{remark}[{[write]} The conjecture and the claim, quoted; what is proved, 6~October 2026]\label{sqt:rem:cmregime}
(a) \emph{The conjecture.} Conjecture~1 of \cite{CM} (arXiv:math/0703600v2,
p.~4) reads: ``Consider a 4-tuple of positive integers $m,s,n,t$ such that
$ms=nt$. Define $\Delta(m,s;n,t)$ by'' \eqref{sqt:eq:CMdelta}. ``Then
$0<\Delta(m,s;n,t)<2$.'' Canfield and McKay record it as proved ``(a) for
$m=n\le9$, using the exact values'' of the Ehrhart polynomials; ``(b) for
sufficiently large $m,n$ when $st=o\bigl((mn)^{1/5}\bigr)$'', using
Greenhill--McKay's asymptotics \cite{GM}; and (c) for several thousand
tuples with $m,n\le30$.

(b) \emph{The claim, as delivered.} The abstract said the density extension
``proves an eventual square case of the Canfield--McKay correction-factor
conjecture''; this section was titled ``An eventual square case of a
correction factor conjecture''; and the sentence after
Corollary~\ref{sqt:cor:delta} began ``This is an eventual
compact-positive-density \emph{square} case of the conjecture, with no
explicit starting index.'' The write narrowed all three on 6~October 2026
(this is the first wording).

(c) \emph{What is proved.} Corollary~\ref{sqt:cor:delta} gives: for every
compact $I\subset(0,\infty)$ there is $n_0(I)$ such that
$0<\Delta(n,s;n,s)<2$ for all $n\ge n_0(I)$ and all integers $s$ with
$s/n\in I$. The threshold $n_0(I)$ comes from the implied constants of
Theorem~\ref{sqt:thm:density} and of the imported theorems and is not
effective, so the corollary verifies the conjecture at no explicitly named
tuple, and it says nothing when $m\ne n$. For each fixed rational density
$\lambda$ it covers the tuples $(n,\lambda n;n,\lambda n)$ with $n$ large,
but not uniformly as $\lambda\to0$ or $\lambda\to\infty$. Square tuples
covered neither by the corollary nor by Canfield and McKay's cases~(a)
and~(b) (their case~(c) is a finite set with $m,n\le30$ that they do not
list) therefore remain for every choice of the $I$'s: sizes $10\le
n<n_0(I)$ within the regime, and, for large $n$, densities $s/n\to0$ with $s$
not $o(n^{1/5})$ (at $m=n$, $s=t$, case~(b) needs $s^2=o(n^{2/5})$), and
$s/n\to\infty$. The corollary settles one regime of the conjecture: not the
conjecture, and not its square case.

(d) \emph{Exact counts.} From the OEIS values (Remark~\ref{sqt:rem:oeis}) and
the definition~\eqref{sqt:eq:CMdelta}, evaluated in 60-digit floating
arithmetic (not interval arithmetic):
\begin{center}\footnotesize\setlength{\tabcolsep}{3pt}
\begin{tabular}{@{}r*{13}{c}@{}}
\toprule
$n$ & 1&2&3&4&5&6&7&8&9&10&11&12&13\\
\midrule
$\Delta(n,n;n,n)$ & 1&1.416&1.291&1.070&0.950&0.869&0.805&0.751&0.704&0.663&0.627&0.594&0.565\\
$n\Delta$ & 1&2.83&3.87&4.28&4.75&5.21&5.63&6.01&6.33&6.63&6.89&7.13&7.35\\
\bottomrule
\end{tabular}
\end{center}
($\Delta(1,1;1,1)=1$ exactly.) Every value lies in $(0,2)$, at distance at
least $0.56$ from both ends; the cases $n\le9$ are within Canfield and
McKay's case~(a). Corollary~\ref{sqt:cor:delta} predicts $n\Delta\to12$.
The products $n\Delta$ increase through $n=13$ and are still far from $12$;
the ratio $\Delta/(12/n)$ rises from $0.36$ at $n=4$ to $0.61$ at $n=13$.
Richardson extrapolation of order $k$ (the model
$n\Delta=c_0+c_1/n+\cdots+c_k/n^k$ fitted exactly at $n=13-k,\ldots,13$)
gives $c_0\approx9.94$, $10.98$, $11.43$, $11.66$, $11.78$ for $k=1,\ldots,5$.
This is consistent with the limit $12$, and it is not a proof of anything;
in particular it bounds no $n_0(I)$.

(e) \emph{The sparse end, formally.} As $\lambda\to0$ the coefficient
in~\eqref{sqt:eq:deltaasymp} behaves like $5/(3un)\sim5/(3\lambda n)=5/(3s)$,
which is the square case $s=t$ of the sparse asymptotic
$\Delta\sim5(s+t)/(6st)$ attributed above to Greenhill--McKay. The two
regimes are thus formally compatible; substituting $\lambda\to0$ into a
fixed-density formula proves nothing about the transition (the source's
Question~4 in Section~\ref{sqt:sec:outlook}; Question~\ref{sqt:q:regime}).
\end{remark}
\begin{writenote}[Independent check of the write, 6~October 2026]
An adversarial check made by the intake after the write (commit
\file{d63de7211}) re-derived the statements marked [write] with its own
code, after fetching again the OEIS entry A110058 and its b-file (revision
and data as quoted), Canfield--McKay's arXiv:math/0703600v2 and Isaev's
arXiv:2508.16952v2. \emph{Lemma~\ref{sqt:lem:sequential}.} The subsequence
argument was re-read and holds; against the rendered preprint, the region~(3.1)
(bounds $(1+\lambda)^{-1}(mn)^{-1/2+2\varepsilon}$,
$(1+\lambda)^{-1}n^{-1/2+\varepsilon}$, $(1+\lambda)^{-1}m^{-1/2+\varepsilon}$
on open half circles), Theorem~3 and the convention of p.~4 are as the
lemma uses them, and~(4.1) at $m=n$ is exactly $I_{0,\lambda}$; the left side
of~(1.1) at $m=n$ is $\tfrac23(4+1/u)$, equal to $3$ at $\lambda=1$ (SymPy).
The paragraph after the lemma matches the two places where the minor-arc
integral enters the proof of Lemma~\ref{sqt:lem:transfer}.
\emph{Remark~\ref{sqt:rem:uniformity}.} Every quotation and every bound
(the factors $\lambda(1+\lambda)^{-1}>(\log n)^{-1}$, $A=O(\log n)$,
$N=\lceil6000(1+\lambda)\rceil$, $\delta=2\pi/N$, the three pieces with
$O(e^{-n^{1-\varepsilon}})I_0$ and $e^{-n^\varepsilon}I_0$, and the omitted
proofs of Lemmas~3 and~4) was found on pp.~15--20; Isaev's Theorem~1.2 and
the inequality $\overline\Delta_V(f,X)\le\Delta_V(f)$ (``Clearly'') were
found as Section~\ref{sqt:sec:provenance} says.
\emph{Remark~\ref{sqt:rem:cmregime}.} The quotations of Conjecture~1 and of
the record of proved cases are verbatim, and the three first wordings are
those of the delivered files. From the fresh b-file, at $80$ digits, all
thirteen values of $\Delta(n,n;n,n)$ and $n\Delta$ in the table, the margin
($\min=0.5651$ at $n=13$), the ratios $0.36$ and $0.61$, both
monotonicities of Question~\ref{sqt:q:monotone} and the Richardson values
$9.94$, $10.98$, $11.43$, $11.66$, $11.78$ were reproduced;
$P_2(1/u)-g_2(u)-\tfrac56=\tfrac{67}{12}+\tfrac5{6u}$ and the expansion of
$(n-1)\log(1+1/n)$ were rechecked in SymPy.
\emph{Remark~\ref{sqt:rem:oeis}.} Revision~\#30, its name, data, link and
Pak's formula are as quoted, Pak's main terms are exactly
$\log L_n-\tfrac12\log(4\pi)-\tfrac14$ (SymPy), and
$\tfrac12\log(4\pi)+\tfrac14=1.515512123\ldots$.
\emph{Remark~\ref{sqt:rem:transseries}.} Parts~(1) and~(3) hold (the master
equation is exact and the triangular recurrence gives $d,e,f,g$). One
statement was too strong: (2) said that the leading balance is not of the
monomial--logarithmic type of \file{plt:thm:lw-template}, but after
$G=\sqrt{F_2/a}$ the reversion is a formal instance of that theorem, as for
the sibling reports \file{a089479-fixed-permanent-matrices} and
\file{a222959-zero-slope-matrices}; corrected with a dated note keeping the
first wording, and the instance added to~(3). \emph{Provenance.} The archive
(669{,}911 bytes, the SHA-256 quoted, 35 files, 34 manifest entries, 980
delivered lines, 24 pages) and the counts of 97 labels and 87 references
were confirmed. For the record, the intake's Richardson values $-1.47$ and
$6.5$ for $D_1$, $D_2$ match fits of order two or three at the top of the
range; the exact fit through all eight values $6\le n\le13$ gives
$-1.4964$ and $6.9139$. No other defect was found.
\end{writenote}

\section{Finite diagnostics and reproducibility}\label{sec:computation}
\subsection{Small exact counts}
A direct dynamic program places one full labelled row at each step. If the residual column capacities are $r_1,\ldots,r_n$, enumerate all vectors $(x_1,\ldots,x_n)$ with $0\leq x_j\leq r_j$ and $\sum_jx_j=n$, then recurse on $(r_j-x_j)_j$. Sorting the residual vector is only a memoization symmetry: distinct labelled row choices remain distinct summands even when they lead to the same sorted state. There is no quotient by column permutations.

The public counter verifies $a_0,\ldots,a_5=1,1,3,55,10147,22069251$. A second unsorted labelled recursion checks $n\leq3$. These values agree with the OEIS fixture \cite{oeis}. Larger tabulated values provide useful diagnostic scale, but they do not establish a remainder constant. For example, let
\[
D_1(n)=n\log(a_n/L_n),\qquad
D_2(n)=n^2\left(\log(a_n/L_n)+\frac3{2n}\right).
\]
The theorem predicts limits $-3/2$ and $223/32=6.96875$, respectively. The finite values in Table~\ref{tab:diagnostics} are still noticeably away from these limits.

\begin{table}[!htbp]
\centering
\caption{Rounded diagnostics from exact counts. The values beyond $n=5$ are source data from OEIS, not recomputed by the public counter.}\label{tab:diagnostics}
\begin{tabular}{rrr}
\toprule
$n$ & $D_1(n)$ & $D_2(n)$\\
\midrule
4 & $-0.753234$ & $2.987065$\\
5 & $-0.851512$ & $3.242439$\\
6 & $-0.918353$ & $3.489884$\\
8 & $-1.011884$ & $3.904924$\\
10 & $-1.077107$ & $4.228934$\\
13 & $-1.146200$ & $4.599394$\\
\bottomrule
\end{tabular}
\end{table}

Neither monotonicity of these diagnostics nor a uniform sign for the omitted terms is claimed. Fitting the finite values would be an unreliable substitute for the exact coefficient calculation.

\subsection{What the package checks}
The companion source archive contains this TeX article, the PDF, fixed-range exact Python implementations, four deterministic JSON receipts, source links, a build script, and an integrity manifest. Its computations cover:
\begin{enumerate}
\item All eighteen connected-Wick Laurent polynomials at costs at most three, with exact rational coefficients, labelled graph multiplicities, and degree bounds.
\item Agreement of each complete polynomial with the independent formal-factorization evaluator, plus separate actual-Gaussian checks at $n=1,2$.
\item Exact local geometric coefficients, density-polynomial identities, the Canfield--McKay correction-factor algebra, and four formal inverse-reversion cancellations.
\item Exact small table counts, input bounds, exclusive-output checks, archive validation, and guards that remain active under optimized Python.
\end{enumerate}
The article supplies the analytic proofs. A receipt does not certify an asymptotic remainder, uniform localization, a finite-size inequality, a numerical inverse ceiling, or a priority claim. Lower Laurent terms present in a low-cost cumulant are retained in the receipts, but are not by themselves complete higher-order coefficients: additional higher-cost cumulants would contribute at those orders.

\subsection{Bounded and read-only reproduction}
The validated code range is cost at most three, at most six vertices, degrees three through eight, and at most nine edges. Resource budgets and type checks use explicit exceptions, including under \texttt{-O}; the integer-to-decimal safety cap remains $640$. The package cannot rewrite its expected receipts. Requested output files must be new and outside the source tree, with exclusive creation and symlink/path checks.

Every Python subprocess uses \texttt{-B}. Both TeX format creation and document compilation explicitly disable shell escape. Builds use disposable copies, verify source hashes before and after, and compare regenerated receipts and PDF bytes with the supplied references. The actual ZIP is extracted independently for ordinary and optimized-Python rebuilds, and both regenerated archives, PDFs and receipts must agree. Byte-identical PDF reproduction depends on a compatible TeX and font toolchain, which the build records; mathematical receipts use deterministic exact arithmetic. The archive includes no third-party paper PDFs, private review reports, or private source paths.

The instructions in \texttt{README.md} and \texttt{COMPUTATION.md} describe commands, dependencies, budgets, failure conditions and the distinction between arithmetic verification and analytic proof. The manifest is an integrity record, not a digital signature. These controls are defensive reproducibility measures, not a sandbox for arbitrary hostile archives or concurrent filesystem attacks.

\section{Context and further questions}\label{sec:outlook}
\subsection{What is inherited and what is specialized}
The leading equivalent is established prior work of Canfield--McKay \cite{CM} and is also within Barvinok--Hartigan's smooth-margin theory \cite{BH}. General complex cumulant control is due to Isaev \cite{Isaev}, with earlier complex-martingale antecedents in Isaev--McKay \cite{IM2018}. Connected-Wick identities and high-dimensional power counting are standard tools, explicitly developed in the related enumeration literature \cite{IM,IMZ}. Higher-order saddlepoint corrections for contingency tables also have a substantial earlier history, including the double-saddlepoint approach of Zipunnikov--Booth--Yoshida \cite{ZBY}.

The assertion proved here is the stated geometric square-table specialization with its complete cutoff and gauge reduction, all-fixed-order rational coefficient algorithm, evaluated corrections, and controlled consequences. It is not a claim to a new general saddlepoint or cumulant method. A bounded primary-source review did not locate the explicit coefficient pair or this dense-square specialization, but absence from the sources inspected is not a worldwide priority certificate. The sparse occurrence of the first numerical coefficient is noted in Section~\ref{sec:conjecture}.

A concrete comparison explains why a leading-order Edgeworth expression does not settle the present correction. Specializing the cubic/quartic exponent displayed by Barvinok--Hartigan to these square margins gives, by direct algebra,
\begin{equation}\label{eq:BHcompare}
-\frac\mu2+\nu=B_0(u)-\frac1n+
\frac1{n^2}\left(\frac7{12}+\frac1{12u}\right).
\end{equation}
At $\lambda=1$, this is $1/4-1/n+5/(8n^2)$. Their theorem uses this expression with relative error tending to zero; it does not assert that its lower-order displayed terms are the true higher-order counting coefficients. The additional connected cumulants account for the difference at the precision sought here. Equation~\eqref{eq:BHcompare} is a specialization of their formula, not a correction to a claim in their theorem.

Recent binary or bounded-entry enumeration results are likewise not interchangeable with geometric tables. For example, admixed arrays in Aw \cite{Aw} and simple multipartite hypergraphs in Isaev--Makai--McKay \cite{IMM} concern different underlying models. Their statements do not supply the analytic specialization above.

\subsection{Questions opened by the expansion}
\begin{enumerate}
\item \textbf{Explicit finite-size bounds.} Can the localization and cumulant constants be tracked sharply enough to certify a useful numerical remainder and a starting index? This would turn the inverse two-ceiling statement and the eventual correction-factor inequality into finite-size certificates.
\item \textbf{Higher coefficients and efficient enumeration.} The finite algorithm proves computability but grows rapidly. Structural reductions of colored connected graphs, and symbolic use of the density polynomials, may make several additional orders practical. The independent moment-factorization route offers a useful second implementation.
\item \textbf{Growth and optimal truncation.} What is the growth of $P_j$ with $j$? Fixed-order proofs do not answer whether the series converges, has a Gevrey bound, or admits useful optimal truncation. Identifying exponentially small contributions would require analysis beyond the present Poincar\'e expansion.
\item \textbf{Transition regimes.} The compact-density constants degenerate as density tends to zero or infinity. A parameter-dependent analysis could connect the positive-density result with sparse table asymptotics or growing-margin regimes. Substitution into a fixed-density formula is insufficient to justify those limits.
\item \textbf{Rectangular and nonconstant margins.} The null direction remains, but the whitening map, covariance structure and coefficient complexity change. Extending the full argument to controlled aspect ratios or smooth nonconstant margins requires fresh uniform bounds rather than an automatic appeal to the square proof.
\item \textbf{The correction-factor conjecture at all sizes.} Corollary~\ref{cor:delta} identifies a positive eventual correction in a dense square regime. Proving the conjectured strict inequalities at every finite size, or uniformly across changing densities and aspect ratios, remains a different problem.
\end{enumerate}
These questions preserve the distinction between the theorem already proved, its computational checks, and stronger conclusions that would need additional analysis.

The theorem numbering for Canfield--McKay and the complex cumulant result refers to the linked preprint versions. The source package links to primary papers but does not redistribute their PDFs.

\renewcommand{\refname}{References and source versions}
\phantomsection\addcontentsline{toc}{section}{References and source versions}
\begin{thebibliography}{99}
\bibitem{oeis}
OEIS Foundation Inc., \emph{The On-Line Encyclopedia of Integer Sequences}, entry A110058, contributed by Brendan McKay.
\url{https://oeis.org/A110058}; numerical data at
\url{https://oeis.org/A110058/b110058.txt}. Accessed 5 October 2026.

\bibitem{CM}
E. R. Canfield and B. D. McKay,
\emph{Asymptotic enumeration of integer matrices with large equal row and column sums},
Combinatorica \textbf{30} (2010), 655--680.
\url{https://doi.org/10.1007/s00493-010-2426-1}.
Preprint title: \emph{Asymptotic Enumeration of Integer Matrices with Constant Row and Column Sums},
arXiv:math/0703600v2, 12 June 2009.
\url{https://arxiv.org/abs/math/0703600v2}.

\bibitem{BH}
A. Barvinok and J. A. Hartigan,
\emph{An asymptotic formula for the number of non-negative integer matrices with prescribed row and column sums},
arXiv:0910.2477v2, 5 April 2010.
\url{https://arxiv.org/abs/0910.2477v2}.

\bibitem{Isaev}
M. Isaev,
\emph{A tail bound for cumulant series for complex functions of independent random variables},
arXiv:2508.16952v2, 28 August 2025.
\url{https://arxiv.org/abs/2508.16952v2}.

\bibitem{IM}
M. Isaev and B. D. McKay,
\emph{Asymptotic enumeration of graph factors by cumulant expansion},
arXiv:2508.18731v1, 26 August 2025.
\url{https://arxiv.org/abs/2508.18731v1}.

\bibitem{IMZ}
M. Isaev, B. D. McKay and R.-R. Zhang,
\emph{Cumulant expansion for counting Eulerian orientations},
Journal of Combinatorial Theory, Series B \textbf{172} (2025), 263--314.
\url{https://doi.org/10.1016/j.jctb.2025.01.002};
\url{https://arxiv.org/abs/2309.15473}.

\bibitem{IM2018}
M. Isaev and B. D. McKay,
\emph{Complex martingales and asymptotic enumeration},
Random Structures \& Algorithms \textbf{52} (2018), 617--661.
\url{https://doi.org/10.1002/rsa.20754};
\url{https://arxiv.org/abs/1604.08305v4}.

\bibitem{GM}
C. Greenhill and B. D. McKay,
\emph{Asymptotic enumeration of sparse nonnegative integer matrices with specified row and column sums},
Advances in Applied Mathematics \textbf{41} (2008), 459--481.
Corrected author version:
\url{https://web.maths.unsw.edu.au/~csg/papers/GMinteger-update.pdf}.

\bibitem{ZBY}
V. Zipunnikov, J. G. Booth and R. Yoshida,
\emph{Counting tables using the double-saddlepoint approximation},
Journal of Computational and Graphical Statistics \textbf{18}(4) (2009).
\url{https://doi.org/10.1198/jcgs.2009.07154}.
Author manuscript:
\url{https://faculty.nps.edu/ryoshida/pdf/jcgs-unblinded.pdf}.

\bibitem{Aw}
A. J. Aw,
\emph{Asymptotic enumeration of admixed arrays and a different independence heuristic},
arXiv:2604.08857v1, 10 April 2026.
\url{https://arxiv.org/abs/2604.08857v1}.

\bibitem{IMM}
M. Isaev, T. Makai and B. D. McKay,
\emph{Enumeration of regular multipartite hypergraphs},
arXiv:2502.08046v2, 21 September 2026.
\url{https://arxiv.org/abs/2502.08046v2}.
\end{thebibliography}

\end{document}
